\documentclass[11pt]{article}

\usepackage[margin=1.15in]{geometry}
\usepackage{amsmath}
\usepackage{amssymb}
\usepackage{amsthm}
\usepackage{booktabs}
\usepackage{longtable}
\usepackage{url}
\usepackage[hidelinks]{hyperref}

\newtheorem{theorem}{Theorem}[section]
\newtheorem{proposition}[theorem]{Proposition}
\newtheorem{lemma}[theorem]{Lemma}
\newtheorem{corollary}[theorem]{Corollary}
\newtheorem{conjecture}[theorem]{Conjecture}
\theoremstyle{definition}
\newtheorem{definition}[theorem]{Definition}
\newtheorem{remark}[theorem]{Remark}

\newcommand{\Phik}{\Phi(K)}
\newcommand{\art}[1]{{\texttt{\hyphenchar\font=`\-\relax #1}}}

\title{Digesting the Four Colour Theorem:\\
schema-robustness of the ring-$14$ obstruction,\\
a minimum-interior theorem, and a coloring-mass law}

\author{%
  (draft --- author list to be completed)%
}

\date{Draft of \today}

\begin{document}
\sloppy
\maketitle

\begin{abstract}
The three machine proofs of the Four Colour Theorem (Appel--Haken 1977;
Robertson--Sanders--Seymour--Thomas 1997; and the 2026 near-linear-time proof of
\texttt{arXiv:2603.24880}) are accepted but opaque: what they establish beyond the
theorem itself is largely unexamined. We report four verified structural facts
obtained by re-deriving these proofs computationally from their published
mathematics rather than by running their released code.

\emph{First and principally}, we show that the failure of the $2026$ discharging
schema over the ring-$\le 14$ sub-pool is a property of the \emph{rule shapes}, not
of the published rule amounts. Fixing the $84$ discharging rule shapes of
\texttt{arXiv:2603.24880} and restricting the published $8{,}200$-configuration
D-reducible pool to its $5{,}895$ members of ring size $\le 14$, we certify that
\emph{no} nonnegative amount vector $x\in\mathbb{R}^{84}_{\ge 0}$ closes the
discharging argument --- ``certify'' in the sense of \S\ref{sec:notclaim}: an
exact-rational Farkas certificate over a row set that is necessary under two
stated hypotheses, not an unconditional proof. The proof is a $12$-row irreducible Farkas certificate,
verified over exact rationals by a solver-free checker that re-derives every row
from its source. The same machinery over the full pool is feasible and the
published amounts satisfy every generated row, so the encoding is not vacuous.
Two hypotheses are load-bearing and stated as such: nonnegativity of the amounts,
and our reconstruction of the degree-$5$/$6$ half of the discharge, which the 2026
paper discharges in a single sentence by citation. This bears directly on a question
Steinberger \cite{steinberger2010} states he has ``no opinion'' about, namely whether
a D-only unavoidable set of ring size $\le 14$ exists; we settle it for this pool and
these rule shapes, and no further.

\emph{Second}, we determine $f(r)$, the minimum number of interior vertices of a
D-reducible configuration of ring size $r$, for $r=8,\dots,13$: $f = 5,5,6,7,7,8$.
The bounds are exhaustive (over $13.2$ million plantri-generated near-triangulations,
$10{,}319$ of which are RSST-legal configurations below threshold, none D-reducible)
and tight (explicit catalog witnesses). This is the first quantitative extension of
Birkhoff's 1913 result that no reducible configuration has ring size $\le 5$.

\emph{Third}, we report an independent re-verification of the published 4CT
configuration corpus by a reducibility checker written from the RSST paper rather
than ported from its code --- the RSST $633$ ($249$ D-reducible, $384$ C-reducible),
Steinberger's $2{,}822$, and $7{,}697$ of the 2026 pool, all with computed
$|\Phi(K)|$ matching the published headers exactly --- together with an end-to-end
reproduction of the 2026 C++ pipeline. Along the way we replaced a legacy C oracle
that aborts partway through Steinberger's set, and identified $430$ configurations
of the 2026 pool that no stage of the pipeline uses, re-verifying the whole pipeline
without them.

\emph{Fourth}, we develop a ``coloring-mass'' framework: exact wheel closed form
$a = (2^r+2(-1)^r)/6$, an improved cap $a \le 11\cdot 2^{r+k-7}$ tight at every
wheel, a proof of $\gamma=3/2$ at $k=2$, and an \emph{obstruction theorem} showing
that no proof factoring through $a \le P(S,4)/24$ can reach the empirically sharp
constant $4/3$.

We are explicit about what did not work. A neural component produced a modest
boundary-classification result and a feature extraction that failed on one of four
folds; it did not discover any theorem here. Three mined sufficient conditions for
reducibility were killed by adversarial mutation search, which produced $236$ false
positives, $215$ of them confirmed by an independent implementation; a natural
LP formulation was killed by its own negative control; the naive threshold law was
killed by the lemma harness. We position the work
relative to the Tao/Buzzard ``digestion'' agenda: the contribution is tooling and
verified structural facts that make an opaque machine proof more comprehensible,
not an automated discovery.
\end{abstract}

\tableofcontents

%======================================================================
\section{Introduction}
%======================================================================

\subsection{The digestion problem}

The Four Colour Theorem has been proved three times by machine
\cite{appel1977every,appel1977every2,rsst1997,steinberger2010,nl4ct2026} and once in
a proof assistant \cite{gonthier2008}. In every case the proof consists of an
\emph{unavoidable set} of configurations, each verified \emph{reducible} by a
finite computation, together with a \emph{discharging} argument certifying
unavoidability. The certificates are enormous --- $633$ configurations and
$32$ discharging rules for RSST \cite{rsst1997}; $2{,}822$ configurations and $42$
rules for Steinberger \cite{steinberger2010}; $8{,}200$ configurations and $84$
rules for the 2026 near-linear-time proof \cite{nl4ct2026} --- and are essentially
never inspected for structure. Correctness has been established; comprehension has
not.

This paper reports on a programme of \emph{computational digestion}: re-deriving
these proofs from their published mathematics with independent implementations, and
then mining the resulting data for statements that a human can hold in mind. The
framing is Tao's and Buzzard's \cite{tao2025machine,buzzard2024}: a machine proof
becomes mathematics when it yields transferable structure, and the bottleneck is
tooling rather than raw computation. The unit of progress we adopt is a
\emph{verified structural fact about the proof objects} --- a theorem about the
class of reducible configurations, or an impossibility statement about a
discharging schema --- rather than a shorter proof of the theorem itself.

Throughout, we distinguish sharply between three epistemic grades, and we label
every claim in the paper with one:
\begin{itemize}
  \item \textbf{Proven}: a mathematical argument, in some cases with a finite
        exhaustive case analysis whose completeness argument is given.
  \item \textbf{Certified}: a finite computation whose output is an independently
        checkable certificate (e.g.\ a Farkas certificate verified in exact
        rationals by a program that does not trust the certificate's own numbers).
  \item \textbf{Corpus-supported}: a statement that survived machine testing over
        our full labelled corpus of $59{,}142$ configurations. This is \emph{not} a
        proof, and \S\ref{sec:negative} records what happened the one time we
        treated it as one.
\end{itemize}

\subsection{The four contributions}

\paragraph{(C1) Schema-robustness of the ring-$14$ obstruction (\S\ref{sec:lp}).}
Steinberger \cite{steinberger2010} produced the first unavoidable set of
\emph{D-reducible} configurations (no contracts), at the cost of ring size up to
$16$ and $2{,}822$ configurations, and wrote that he had ``no opinion'' whether an
unavoidable set of D-reducible configurations with ring size at most $14$ exists.
We attack a sharply scoped version of this question. Fix the $84$ rule shapes of
\cite{nl4ct2026} and the sub-pool $P_{14}$ of the published $8{,}200$-configuration
D-reducible pool consisting of the $5{,}895$ configurations of ring size $\le 14$.
Then there is no $x \in \mathbb{R}^{84}_{\ge 0}$ under which the discharging
argument closes over $P_{14}$; a $12$-row irreducible Farkas certificate proves it,
verified in exact rational arithmetic without a solver. The contradiction lives
entirely at the degree-$5$/degree-$7$ interface, and $65$ of the $84$ rule shapes
are irrelevant to it. Because the published amounts satisfy the analogous full-pool
system with zero violations, the result distinguishes the two pools rather than
being vacuously infeasible.

\paragraph{(C2) The function $f(r)$ (\S\ref{sec:fr}).} Birkhoff \cite{birkhoff1913}
proved that a reducible configuration has ring size at least $6$. No quantitative
refinement --- a lower bound on the \emph{interior} size as a function of the ring
size --- appears to have been published. We prove
\[
  f(8),\dots,f(13) \;=\; 5,\,5,\,6,\,7,\,7,\,8,
\]
exhaustively and tightly, and note the nuance that the corresponding minimum over
\emph{C}-reducible configurations is strictly smaller at rings $8$ and $11$.

\paragraph{(C3) Independent verification of the corpus (\S\ref{sec:verify}).} We
re-verified the published catalogs with a checker implemented from the RSST paper
\cite{rsst-reduce} rather than ported from \texttt{reduce.c}, reproduced the 2026
pipeline end to end (four stages, $4$h$38$m wall, asserts live, all counts matching
the paper), replaced a legacy C oracle where it aborts, and identified $430$
configurations of the 2026 pool that no stage of the pipeline uses.

\paragraph{(C4) The coloring-mass law (\S\ref{sec:mass}).} Writing $a=|\Phi(K)|$ for
the classical RSST count of ring colorings extending to the free completion, we
prove the exact wheel value $a=(2^r+2(-1)^r)/6$ with its equality case, the cap
$a \le 11\cdot 2^{r+k-7}$ (tight at every wheel), and $\gamma=3/2$ at $k=2$. We
give an \emph{obstruction theorem}: an explicit $10$-vertex configuration shows that
no bound factoring through $a \le P(S,4)/24$ can prove the empirically sharp cap
$a \le \frac{2^r+2}{6}(4/3)^{k-1}$, the floor for that route being $61/43$. Paired
with an empirical threshold law, the sharp cap would reproduce the entire $f(r)$
table analytically, to within one interior vertex at every ring.

\subsection{What this paper does not claim}
\label{sec:notclaim}

We state the limitations up front because several of them are load-bearing.

\begin{enumerate}
  \item \textbf{(C1) does not settle Steinberger's question.} It is scoped to
        \emph{these} $84$ rule shapes and \emph{this} pool. Other D-reducible
        configurations of ring $\le 14$ exist outside the published pool, and
        adding or replacing rule shapes is not covered. Removing the fixed-shapes
        caveat requires column generation over planar rule patterns
        (\S\ref{sec:open}).
  \item \textbf{Unblocked $\ne$ realizable.} Our LP rows are generated for every
        local structure the pool fails to block, exactly as the released C++ decides
        which cartwheels to analyse. A structure that is geometrically impossible
        for a reason outside the pool would yield a spurious row. This is
        anti-conservative for an \textsc{infeasible} verdict; it is also precisely
        the standard the published proof itself uses.
  \item \textbf{The degree-$5$/$6$ rows are a reconstruction.} The released code
        never enumerates hubs of degree $\le 6$; \cite{nl4ct2026} discharges that
        case in one sentence by citation to \cite{steinberger2010}, who proved it
        \emph{by machine} at his own amounts. \S\ref{sec:lowdeg} gives our
        reconstruction, its verbatim textual audit against the paper source, and an
        exhaustive refinement check. Three of the $12$ certificate rows depend on it.
  \item \textbf{The sharp cap is a conjecture}, not a theorem, and the threshold
        half of the coloring-mass law is corpus-supported only.
  \item \textbf{No machine-learning discovery is claimed.} \S\ref{sec:ml} reports
        the neural work honestly, including a replication finding that the same
        architecture at the same accuracy learned a \emph{different} mechanism than
        a published result of the same shape.
  \item \textbf{Schema-incomparability.} Configuration counts across schemas are not
        directly comparable: RSST's $633$ uses an axle-based discharging schema with
        C-reducibility and contracts at ring $\le 14$; the 2026 schema is
        wheel/cartwheel-based with D-reducibility only. Any claim about set size
        must be schema-qualified.
\end{enumerate}

\subsection{Related work}

\paragraph{On $f(r)$.} Birkhoff's ring-size lower bound \cite{birkhoff1913} is, as far
as we can determine, the only published unconditional structural lower bound on the
size of a reducible configuration. The literature contains extensive
\emph{catalogs} --- from which per-ring interior minima can be read off, and which we
use as our tightness witnesses --- but a catalog minimum is an upper bound on $f(r)$,
not a lower bound; the exhaustive negative side of Theorem~\ref{thm:fr} is what is
new. We have found no prior claim of a per-ring interior minimum beyond
Birkhoff's ring $\ge6$, and no prior exhaustive enumeration of the small-interior
configuration classes. The nuance that $C$-reducibility admits smaller interiors
(\S\ref{sec:cvsd}) is implicit in the RSST catalog but, to our knowledge, has not
been stated.

\paragraph{On lower bounds for unavoidable sets.} Bounds on the \emph{number} of
configurations needed are essentially absent from the literature; the historical
sequence $\sim1834$ \cite{appel1977every2} $\to 633$ \cite{rsst1997} $\to 2{,}822$
D-only \cite{steinberger2010} $\to 8{,}200$ D-only with a near-linear colouring
algorithm \cite{nl4ct2026} records only what has been achieved. Theorem~\ref{thm:lp}
is not a bound on set size; it is an impossibility statement about a specific
(schema, rule-shape set, pool) triple, and we know of no comparable prior result.
The closest antecedent in spirit is Steinberger's own remark that he has no opinion
about the ring-$14$ question.

\paragraph{On digestion.} The programme of making machine proofs comprehensible ---
rather than merely correct --- is articulated by Tao \cite{tao2025machine} and
Buzzard \cite{buzzard2024}, and Gonthier's Coq formalisation \cite{gonthier2008}
settles the correctness half for the 4CT specifically. Our contribution is on the
comprehension side and is deliberately modest in ambition: verified structural facts
about the proof objects, plus the tooling that produced them.

\paragraph{On machine learning for mathematical discovery.} The template we followed
is that of Davies et al.\ \cite{davies2021,davies2021signature}: train a model, probe
it, extract a human-statable conjecture. \S\ref{sec:ml} reports our attempt and its
partial failure, together with a replication study bearing on the reliability of the
``extract a mechanism'' step in general. We regard the honest reporting of that
attempt as part of the contribution, not as an appendix to it.

\subsection{Reproducibility}

Every numerical claim in this paper is annotated with the repository artifact that
produced it. Appendix~\ref{sec:artifacts} is a claim-to-artifact table giving, for
each claim, the file path and what it contains. All third-party inputs (the four
arXiv e-prints, the \texttt{plantri} generator, the released rule and configuration
files, the released C++ checker) are pinned in-repository with recorded SHA-256
digests.

%======================================================================
\section{Preliminaries}
\label{sec:prelim}
%======================================================================

We follow \cite{rsst-reduce,rsst1997} throughout.

\subsection{Configurations and free completions}

\begin{definition}[Configuration]
A \emph{configuration} $K$ of ring size $r$ with $k$ interior vertices is given by
its \emph{free completion} $S=S(K)$: a simple plane near-triangulation of the closed
disk on $n=r+k$ vertices such that
\begin{enumerate}
  \item the outer face is bounded by an induced (chordless) cycle $R$, the
        \emph{ring}, of length $r$, each of whose vertices has degree $\ge 3$;
  \item the \emph{interior} (the $k$ vertices strictly inside) is nonempty and
        connected, and every interior vertex has degree $\ge 5$;
  \item every bounded face is a triangle.
\end{enumerate}
\end{definition}

RSST's reference parser \texttt{ReadConf} in \texttt{reduce.c} enforces two further
conditions on any file it accepts, and we enforce both:
\begin{enumerate}
  \setcounter{enumi}{3}
  \item a \emph{labelling convention}: each ring vertex's neighbour list, written in
        rotation order, starts at ring-neighbour $i+1$ and ends at ring-neighbour
        $i-1$, with every entry strictly between them interior;
  \item the degree identity $\sum_v \deg v = 6(n-1)-2r$ (automatic given 1--3);
  \item \textbf{each interior vertex meets the ring in at most two separate
        contiguous arcs.}
\end{enumerate}
Condition 6 is genuine structure, not bookkeeping: it is not implied by 1--3, and it
removes $33\%$--$77\%$ of the near-triangulations satisfying 1--3 at the sizes we
enumerate (\S\ref{sec:fr-enum}). Condition 4 is a serialization convention only; a
cyclic neighbour list denotes the same rotation system whatever entry it is written
from. Discovering that distinction cost us a real bug (\S\ref{sec:methods}).

\subsection{Reducibility}

Let $\Phi(K)$ be the set of colourings of $E(R)$ in three colours, taken up to the
$S_3$ action on colours, that extend to a \emph{tri-colouring} of $S$ --- an edge
$3$-colouring in which the three edges of every internal face receive three distinct
colours. Set $a = a(K) = |\Phi(K)|$; this is the classical header field of an RSST
\texttt{.conf} record. $K$ is \emph{D-reducible} if the Kempe-closure operator on
subsets of $\Phi(K)$, iterated to its fixed point starting from $\Phi(K)$, reaches
$\emptyset$; equivalently, no consistent set survives. $K$ is \emph{C-reducible} if
the same holds after contracting a nonempty \emph{contract} $X$ of edges. C-reducibility
is strictly weaker: \S\ref{sec:cvsd} exhibits catalog configurations that are
C-reducible but not D-reducible with fewer interior vertices than $f(r)$.

\begin{remark}
Everything in \S\ref{sec:fr} is scoped to D-reducibility. The $2026$ pool
\cite{nl4ct2026} and Steinberger's set \cite{steinberger2010} are D-only; the RSST
$633$ are not.
\end{remark}

\subsection{Discharging and the three schemas}

A discharging argument assigns an initial charge to each vertex of a hypothetical
minimum counterexample, moves charge along edges according to a finite set of local
\emph{rules}, and derives a contradiction from the fact that total charge is
positive while every vertex ends non-positive. RSST use \emph{axles}; Steinberger
and \cite{nl4ct2026} use \emph{cartwheels} (a hub with its first and second
neighbourhoods, degrees given as intervals). The 2026 paper assigns
$T_0(v) := 10\,(6-d(v))$, so $\sum_v T_0(v) = 120$ by Euler
(\cite{nl4ct2026} source line 602 and Lemma \texttt{lem:120}); a rule $R$ sends
$r(R)\in\{1,2\}$ units from $s(R)$ to $t(R)$; and the final charge is
$T(u) = T_0(u) + \sum_{v\sim u}(\phi(v,u)-\phi(u,v))$.

Crucially, \emph{survivors are permitted}: the charge argument need not kill every
wheel, because three gluing lemmas mop up the survivors. Encoding ``the discharging
must kill every wheel'' is therefore the \emph{wrong} condition, and
\S\ref{sec:strongenc} reports the negative control that demonstrates it: that
encoding is infeasible even for the full $8{,}200$-configuration pool, i.e.\ it would
``disprove'' the published theorem.

\subsection{Computational setting}

Our reducibility checker \texttt{fourcolor.reduce.check} is implemented from
\cite{rsst-reduce}'s mathematics (the \texttt{reduce.tex} write-up) rather than
ported from \texttt{reduce.c}; the 1995 C program is compiled separately as an
oracle (\texttt{build/reduce\_rsst}) purely for differential testing. Exhaustive
generation uses \texttt{plantri} 5.5 \cite{brinkmann2007}, pinned with its
SHA-256 digest, invoked as \texttt{plantri -P$r$ -c3 -a $n$}; with \texttt{-P} and
no \texttt{-m}, \texttt{-c3} defaults to \texttt{-c3m3}, i.e.\ no chords and no
outer-face vertex of degree $2$, exactly matching condition 1. \texttt{plantri} is a
canonical-construction-path generator producing exactly one representative per
isomorphism class, which is what makes ``sweep every $n$ below threshold'' an
exhaustive search rather than a sample.

%======================================================================
\section{Schema-robustness of the ring-$\le 14$ obstruction}
\label{sec:lp}
%======================================================================

\subsection{The question and the exact claim}

Steinberger's unavoidable set \cite{steinberger2010} is D-only but reaches ring size
$16$; RSST's $633$ stay at ring $\le 14$ but use contracts. Steinberger states the
question and declines to conjecture
(\art{third\_party/arxiv-0905.0043/src/4c.tex}, line 137):
\begin{quote}\small
``Appel and Haken used only configurations of ring-size at most $14$ \dots\ The
unavoidable set of Robertson et al.\ also has only configurations of ring-size $14$
or less. Our unavoidable set uses configurations of ring-size up to $16$ \dots\
\textbf{We were not able to find an unavoidable set using only configurations of
ring-size $14$ or less, but we have no opinion as to whether such a set exists or
not.}''
\end{quote}
The 2026 pool \cite{nl4ct2026} is D-only with $8{,}200$ configurations of ring size
up to $18$, of which $5{,}895$ have ring size $\le 14$.

A preliminary experiment (\S\ref{sec:s1}) ran the released 2026 pipeline with the
pool restricted to those $5{,}895$: it fails, with $833$ of $16{,}157$ per-wheel jobs
aborting on an assertion. That is a statement about the \emph{published amounts}
$x_0$. The result of this section is the much stronger statement that no re-tuning
of the amounts can rescue it.

\begin{theorem}[Certified]
\label{thm:lp}
Fix the $84$ discharging rule shapes of \cite{nl4ct2026}
(\art{third\_party/discharging-rules/R}). Let $P_{14}$ be the sub-pool of the
published $8{,}200$-configuration D-reducible pool consisting of the $5{,}895$
configurations of ring size $\le 14$. Then there is no vector of nonnegative rule
amounts $x \in \mathbb{R}^{84}_{\ge 0}$ satisfying the local charge conditions that
the $2026$ schema imposes over $P_{14}$: i.e.\ no $x\ge 0$ under which every
cartwheel leaf that $P_{14}$ fails to block satisfies $C(L,x)\le 0$ (strictly
$<0$ in the cases of \S\ref{sec:necessity}) and every unblocked degree-$5$ or
degree-$6$ wheel satisfies $\mathrm{LB}(w,x)\le 0$. In particular the discharging
argument of \cite{nl4ct2026} cannot be made to close over $P_{14}$ by re-tuning
its amounts. The hypotheses under which these rows are necessary conditions ---
including that an unblocked local structure need not be realizable --- are stated
in \S\ref{sec:lp-scope}.
\end{theorem}

\begin{remark}[$43$ figures, $84$ rules]
\cite{nl4ct2026} draws $43$ rules in its Figure \texttt{fig:rules}; each is counted
twice, once per reflection, except two that are reflection-symmetric, giving
$43\cdot2-2 = 84$ (source lines 635--636, ``Figure \texttt{fig:rules} shows 84 rules
since precisely two, the first and fourth last, are symmetric under the
reflection''). This matches the $84$ files in
\art{third\_party/discharging-rules/R} one for one, and the paper notes the rule set
is the same as Steinberger's $42$ except that his $12$th and $13$th are merged. Our
$x$ ranges over all $84$ independently, which is weaker than assuming reflection
pairs share an amount, hence the safe direction for \textsc{infeasible}.
\end{remark}

The certificate is $12$ rows; it is stated in full in
Appendix~\ref{sec:certificate}. Since all $12$ rows are non-strict, infeasibility
holds over the reals, hence a fortiori over the integers and over the box
$\{1,2\}^{84}$ that \cite{nl4ct2026} allows itself. Section~\ref{sec:lp-scope}
lists the hypotheses.

\subsection{The empirical precursor: the pipeline on the restricted pool}
\label{sec:s1}

Before any LP, we simply ran the released pipeline with the pool replaced by
$P_{14}$ and the $84$ rule shapes and published amounts $x_0$ left untouched
(\art{tools/s1\_build\_pool.py}, \art{results/p3/runs/steinberger-s1/}). The
excluded $2{,}305$ configurations are $1{,}695$ at ring $15$, $526$ at $16$, $69$
at $17$ and $15$ at $18$.

The schema breaks, and it breaks in a structured way.
\texttt{combine\_rules} yields $681$ non-blocked combined rules rather than the
baseline $671$ --- ten combined rules revived, all ten predicted in advance by a
gap analysis with $0$ unexpected re-blockings. Wheel counts are unchanged at
$d=7,8$ and grow by $4$ and $5$ at $d=9,10$, all nine revivals again predicted ---
so the per-wheel job count is $16{,}148 + 9 = 16{,}157$ rather than the full pool's
$16{,}148$.
Then \textbf{$833$ of $16{,}157$ per-wheel cartwheel jobs abort}:
$780$ on \texttt{assert(C == 0)} ($530$ at $d=7$, $250$ at $d=8$) and $53$ on
\texttt{assert(d == 7 || d == 8)} ($51$ at $d=9$, $2$ at $d=10$). The two failure
kinds are perfectly stratified by hub degree, which is itself informative. The run's
verdict is \textsc{fail} at the \texttt{enum\_cartwheels} stage.

This says the published amounts do not work over $P_{14}$. It does not say no
amounts do --- and the obvious response, ``re-tune $x$'', is exactly what
Theorem~\ref{thm:lp} rules out. The $833$ aborted jobs are also the raw material for
the LP: re-run against an NDEBUG build with the assertions compiled out, they yield
$77{,}185$ additional fully refined leaves that the assert-guarded run never
recorded.

\subsection{Encoding: leaf charge is affine in $x$}

The invariant the released C++ actually enforces
(\art{third\_party/computer-checks/src/cartwheel.cpp}, \texttt{enum\_bad\_cartwheels})
is, on every fully refined cartwheel $L$ surviving \texttt{fix\_in\_rules} and
\texttt{fix\_out\_rules}, the conjunction
\texttt{assert(C==0)}, \texttt{assert(d==7||d==8)},
\texttt{assert(darts\_by\_deg[7]+[8]+[9] > 0)}.
The corresponding necessary condition on $x$ is:

\begin{center}
\begin{tabular}{lll}
\toprule
case & requirement & why \\
\midrule
$d(L)\in\{9,10,11\}$ & $C(L,x) < 0$ & $L$ must be pruned before the degree assert \\
$d(L)\in\{7,8\}$, no spoke of degree $\ge 7$ & $C(L,x) < 0$ & pruned before the spoke assert \\
$d(L)\in\{7,8\}$, some spoke $\ge 7$ & $C(L,x) \le 0$ & $C=0$ survives to the gluing lemmas \\
\bottomrule
\end{tabular}
\end{center}

Note that $C=0$ is \emph{not} imposed: if $C<0$ the leaf is pruned before the assert
is reached, so $C \le 0$ is exactly the necessary-and-sufficient charge condition at
a leaf, weaker than what the C++ observes at $x_0$.

At a leaf all $d$ spokes are fixed, so the per-spoke maximum inside
\texttt{amount\_of\_possible\_charge\_send} never runs, and
\texttt{prune\_by\_non\_associated\_rule} forces the spoke combined rules to be
determined by the leaf graph. Hence
\begin{equation}
\label{eq:leafrow}
  C(L,x) \;=\; 10\,(6-d) \;+\; \sum_{k} (\mathrm{IN}_k - \mathrm{OUT}_k)\, x_k,
\end{equation}
with $\mathrm{IN}_k = \#\{j : \texttt{always\_apply}(\text{dart}_j, R_k)\}$ counting
darts pointing into the hub and $\mathrm{OUT}_k$ the reverse. The form is affine
with integer coefficients: no epigraph variables, no binaries.

\paragraph{Differential validation of \eqref{eq:leafrow}.} All $10{,}094$ published
full-pool bad cartwheels in \art{third\_party/computer-checks/wheels/zero} evaluate
to \emph{exactly} $0$ at $x_0$ under the Python form, and all $10{,}094$ satisfy
$d\in\{7,8\}$ ($9{,}366$ at $d=7$, $728$ at $d=8$) with a spoke of degree $\ge 7$;
i.e.\ the Python row semantics independently reproduce all three C++ assertions on
the published output, with $0$ discrepancies
(\art{tools/lp\_discharge.py -{}-validate}).

\subsection{Why leaves observed at $x_0$ constrain every $x$}

The refinement forest is explored with pruning, so the observed leaf set depends on
$x_0$. The rows are nevertheless valid for all $x \ge 0$:

\begin{lemma}[Monotonicity]
\label{lem:mono}
For every observed leaf $L$ and every $x \ge 0$, the necessary condition
$C(L,x)\le 0$ (resp.\ $<0$ in the strict cases) holds.
\end{lemma}

\begin{proof}[Proof sketch]
(a) The branching operations \texttt{update\_degree\_by\_rule},
\texttt{concrete\_degree\_except\_tail}, \texttt{refine\_always} and
\texttt{refine\_never} only narrow degree intervals and never change rotations, so
branching is $x$-independent. It is \emph{not} pool-independent, since
\texttt{fix\_in\_rules} branches over the pool-dependent non-blocked combined-rule
set --- but it moves in the safe direction: a smaller pool blocks fewer combined
rules, and we verified that the full-pool combined rules are a content-exact subset
of the $P_{14}$ ones ($671 \subset 681$ files), so the full-pool forest is a
sub-forest of the $P_{14}$ forest.
(b) \texttt{blocked\_by\_reducible\_configuration} is $x$-independent and monotone in
the pool; \texttt{prune\_by\_non\_associated\_rule} is independent of both.
(c) \texttt{upper\_bound\_of\_charge} is non-increasing along refinement whenever
$x \ge 0$: narrowing degrees can only add \texttt{always\_apply} matches (out-charge
up --- this is the step needing $x\ge 0$) and only remove candidates surviving
\texttt{never\_apply}.
(d) Hence for any $x\ge 0$, either every ancestor survives and the assert demands
$C(L,x)\le 0$, or some ancestor $A$ has $C(A,x)<0$ and then $C(L,x)\le C(A,x)<0$ by (c).
\end{proof}

An earlier draft asserted flat pool-independence in (a). That was false and was
caught by the adversarial pass (\S\ref{sec:methods}); the corrected statement is
what the argument needs and it holds. Since $P_{14}\subseteq P_{\text{full}}$
(verified: $0$ of the $5{,}895$ names absent from the $8{,}200$; both are symlinks to
the same files), every full-pool leaf row is also valid for the ring-$\le 14$
problem, which is what lets the certificate mix rows of both origins.

\begin{remark}[Direction of the relaxation]
The row set omits the three gluing lemmas and omits every leaf the $x_0$-pruned
search never reached. Both omissions make the encoded feasible set \emph{larger}
than the true one. Therefore \textsc{infeasible} transfers to the true problem,
while \textsc{feasible} is only a non-refutation.
\end{remark}

\subsection{The degree $\le 6$ half --- without which the LP is vacuous}
\label{sec:lowdeg}

The released code enumerates hubs of degree $7$--$11$ only, where the initial charge
$10(6-d)$ is already negative. Taken alone the leaf rows are satisfied by $x=0$:
send nothing and every hub keeps its negative initial charge. We measured exactly
that --- the ``no low-degree'' variants in Table~\ref{tab:lpvariants} are
\textsc{feasible} and vacuously so.

The 2026 paper does not check $d\le 6$ by hand or otherwise. It asserts the
\emph{stronger} $T(v)=0$ and discharges the obligation in one sentence by citation,
inside the proof of its main theorem (source line 1144):
\begin{quote}\small
``\dots we are still missing the case where $v$ has degree at most $6$ and positive
final charge. However, this cannot happen, for if $v$ has degree 5 or 6, then it has
final charge 0 as proved in \cite{steinberger2010}, where exactly the same
discharging rules as here are used.''
\end{quote}
And \cite{steinberger2010} does not prove it by hand either: he proves it by machine,
with one presentation file per hub degree from $5$ to $11$, and explicitly declines
to write the hand argument
(\art{third\_party/arxiv-0905.0043/src/4c.tex}, line 552). So the $d\le 6$ case rests
on a computation at Steinberger's amounts and against Steinberger's unavoidable set.
When the LP re-tunes $x$, that citation no longer applies as proved --- which is
exactly why we impose the necessary condition $T(v)\le 0$ directly as a row rather
than inheriting the conclusion.

For $d\in\{5,6\}$ we add rows built from a \emph{lower} bound on the hub's final
charge, valid for any $x\ge 0$ and at any level of refinement:
\begin{equation}
\label{eq:lowdeg}
  \mathrm{LB}(w,x) \;=\; 10(6-d) \;-\; \sum_i \mathrm{POSS}_i(x) \;+\; \sum_j \mathrm{ALW}_j(x)
  \;\le\; T(v) \;\le\; 0,
\end{equation}
where $\mathrm{POSS}_i$ sums $x_k$ over base rules \emph{not} \texttt{never\_apply}
out of spoke $i$ (an over-count of sends) and $\mathrm{ALW}_j$ sums $x_k$ over base
rules that \texttt{always\_apply} into the hub from spoke $j$ (an under-count of
receipts). Rows are generated over every spoke-degree necklace in $\{5,\dots,9\}^d$
up to rotation that the pool does not block.

\paragraph{Textual audit.} We audited \eqref{eq:lowdeg} quote by quote against the
paper's pinned LaTeX source. Initial charge $10(6-d)$: verbatim, line 602. Total
charge $120$: Lemma \texttt{lem:120}. $T = T_0 - \text{out} + \text{in}$: Definition
\texttt{dfn:finalcharge}. Requirement at $d\in\{5,6\}$: the paper asserts $T(v)=0$,
which is \emph{stronger} than our $T(v)\le 0$, so our rows are conservative. A
degree-$5$/$6$ vertex with $T=0$ is a permitted survivor, matching our non-strict
rows. Rule $R(1)$ decoded from the paper's own figure source
(\art{tikz/rule.tex}, lines 967--977) is a single edge with source of degree exactly
$5$, target of degree $\ge 5$, amount $2$, matching \art{rule001.rule} byte for byte;
so a degree-$5$ vertex sends $2$ along every incident edge unconditionally, and the
row $10 - 5x_{\text{rule001}} \le 0$ is forced. Full audit:
\art{results/steinberger/degree56-verification.md}.

\paragraph{The coarse-cartwheel gap, and its closure.} The low-degree rows are built
on the \emph{coarse} cartwheel (hub plus spoke-degree necklace, second neighbourhood
left at $[5,9]$), and blocking is tested on that coarse object. ``Refining only
strengthens the bound'' does \emph{not} cover the danger, because refining can also
\emph{delete the row}: if a coarse wheel is unblocked but every full refinement is
blocked, then no vertex of a minimum counterexample carries that necklace and the row
is spurious. That is anti-conservative for \textsc{infeasible}, so it was checked
rather than assumed (\art{tools/lowdeg\_refinement\_check.py}).

\emph{No low-degree row is spurious.} Two findings, the second strictly stronger than
needed. (i) Exact enumeration of the certificate's four degree-$5$ necklaces over the
full refinement space $\{5,\dots,9\}^m$, $m=10$--$12$, by compiling the blocking
predicate into a DNF over the open slots and counting models exactly by DFS; the
compiled predicate is differentially checked against the production
\texttt{wheel\_is\_blocked} on $2{,}000$ random refinements per necklace ($0$
disagreements) and exhaustively on all $625$ refinements of a $4$-slot necklace.

\begin{center}
\begin{tabular}{lrrrrr}
\toprule
necklace & slots & refinements & blocked & unblocked & surviving \\
\midrule
\texttt{d5-55767} (IIS) & 10 & 9{,}765{,}625   & 5{,}678{,}597  & 4{,}087{,}028   & 41.9\% \\
\texttt{d5-57577} (IIS) & 11 & 48{,}828{,}125  & 30{,}711{,}497 & 18{,}116{,}628  & 37.1\% \\
\texttt{d5-57677} (IIS) & 12 & 244{,}140{,}625 & 94{,}647{,}089 & 149{,}493{,}536 & 61.2\% \\
\texttt{d5-55677} (19-row support) & 10 & 9{,}765{,}625 & 6{,}350{,}608 & 3{,}415{,}017 & 35.0\% \\
\bottomrule
\end{tabular}
\end{center}

Not a knife-edge: between a third and two-thirds of all refinements of each necklace
remain unblocked. (ii) The gap closes uniformly for all $86$ low-degree rows at once:
\texttt{PseudoConfiguration::representative\_degree} collapses any vertex with degree
upper bound $>8$ to the single value $9$, so the coarse blocking test \emph{is} the
blocking test of the refinement pinning every second neighbour to $9$. Verified
mechanically over the whole sweep: all $580$ unblocked $d=5$ and all $2{,}466$
unblocked $d=6$ necklaces have an unblocked tail-maximised refinement, $0$ exceptions
of $3{,}046$.

\paragraph{Sanity anchor.} At $x_0$, the maximum of $\mathrm{LB}$ over all $3{,}046$
unblocked low-degree wheels is exactly $0$, never positive; it is attained e.g.\ at
$d=5$ with all five spokes of degree $9$, where the row is literally
$10 - 5x_{\text{rule001}} \le 0$ and $x_0$ has $x_{\text{rule001}}=2$. So $x_0$ sits
simultaneously on the boundary of the low-degree rows and of all $3{,}425$ distinct
full-pool leaf rows --- the signature of a tightly optimised scheme, and a strong
consistency check: had we got the $d\le 6$ side wrong in the loose direction, $x_0$
would have violated it. Independently, running the released C++ with
\texttt{-{}-enum\_wheels -d 5} and \texttt{-d 6} returns exactly $580$ and $2{,}466$
unblocked wheels, matching our generation wheel for wheel.

\subsection{Code-necessary versus math-necessary}
\label{sec:necessity}

This is the objection that nearly killed the result. \texttt{upper\_bound\_of\_charge}
counts only \texttt{always\_apply} rules outward, and \texttt{fix\_out\_rules} stops
refining once no rule is \texttt{dominantly\_apply}-but-not-\texttt{always\_apply}.
But \texttt{dominantly\_apply} is strictly stronger than \texttt{has\_intersection},
so a leaf can retain rules that are neither \texttt{always\_apply} nor
\texttt{never\_apply} on an out-dart --- rules that might fire in some realization but
are not counted in \texttt{out\_charge\_sum}. Measured: $11$ of $5{,}292$
(rule, out-dart) pairs across the certificate's leaves are undetermined. For those
leaves $C$ strictly \emph{over}-estimates the hub's true final charge, so
$C(L,x)\le 0$ would be a condition of the released \emph{verification procedure}
rather than of the argument itself.

Two blunt diagnostics confirm this is not hypothetical --- \textbf{both make the
ring-$\le14$ LP feasible} (Table~\ref{tab:lpvariants}, last four rows). Both,
however, are needlessly lossy: one throws rows away, the other weakens rows that did
not need weakening.

\paragraph{The repair.} Pass to the \emph{tail-maximised} leaf: narrow every open
tail $[a,9]$, $a<9$, to $[9,9]$. Since $9$ means ``$9$ or more'' in the cartwheel
abstraction, this restricts attention to the sub-family of realizations in which
every unpinned second neighbour has large degree. If the tail-maximised leaf
(1) has a fully determined OUT side, (2) yields the identical row --- same
coefficients, same constant, same strictness --- and (3) is still unblocked by the
pool, then in those realizations OUT is exact while IN can only exceed the branch's
flag sum, so the true final charge is $\ge C(L,x)$, and $T\le 0$ forces
$C(L,x)\le 0$. Undetermined IN rules are harmless precisely because they push the
true charge up.

\emph{All $12$ certificate rows pass this audit.} The nine degree-$7$ rows satisfy
(1)--(3) individually, including the five whose original leaves had undetermined
out-pairs; the three degree-$5$ rows are mathematically necessary by construction,
being already in lower-bound form. Machine-readable:
\art{results/steinberger/iis\_math\_necessity.json},
\texttt{all\_mathematically\_necessary: true}.

\subsection{Results}

\begin{table}[t]
\centering\small
\caption{LP variants. Rows are deduplicated by (coefficients, constant, strictness).
Reproduce with \art{tools/lp\_discharge.py -{}-all} ($\approx$14 min, $\approx$2 GB peak).}
\label{tab:lpvariants}
\begin{tabular}{lrrrl}
\toprule
variant & rows & strict & $x_0$ viol. & status \\
\midrule
\textbf{control}, full pool, integral, with $d\le6$ & 3{,}511 & 0 & \textbf{0} & \textbf{feasible} \\
control, full pool, real, with $d\le6$ & 3{,}511 & 0 & 0 & feasible \\
control, full pool, integral, no $d\le6$ & 3{,}425 & 0 & 0 & feasible (vacuous) \\
control, full pool, real, no $d\le6$ & 3{,}425 & 0 & 0 & feasible (vacuous) \\
\textbf{main}, ring$\le$14, integral, with $d\le6$ & 48{,}509 & 96 & 9{,}647 & \textbf{infeasible} ($y^\top b=-30$, 19 rows) \\
\textbf{main}, ring$\le$14, real, with $d\le6$ & 48{,}509 & 96 & 9{,}551 & \textbf{infeasible} ($y^\top b=-30$, 19 rows) \\
main, ring$\le$14, integral, no $d\le6$ & 48{,}423 & 96 & 9{,}647 & feasible (vacuous, $x=0$) \\
main, ring$\le$14, real, no $d\le6$ & 48{,}423 & 96 & 9{,}551 & feasible (vacuous, $x=0$) \\
\midrule
control, \texttt{-{}-exact-only} (\S\ref{sec:necessity}) & 2{,}620 & 0 & 0 & feasible \\
main, \texttt{-{}-exact-only} & 42{,}121 & 96 & 8{,}699 & feasible --- \emph{rows dropped, too lossy} \\
control, \texttt{-{}-conservative} & 4{,}321 & 0 & 0 & feasible \\
main, \texttt{-{}-conservative} & 52{,}400 & 96 & 8{,}741 & feasible --- \emph{rows weakened, too lossy} \\
\bottomrule
\end{tabular}
\end{table}

Leaf inputs for the main run: $97{,}860$ fully refined cartwheels yielding $48{,}423$
distinct leaf rows, to which the $86$ distinct low-degree rows are added, giving
$48{,}509$. Composition of the leaves: $10{,}094$ published full-pool leaves; $10{,}581$ leaves
from the $15{,}324$ per-wheel jobs that completed; $77{,}185$ leaves recovered from
the $833$ jobs that aborted on an assert, by re-running them against an NDEBUG build
with the asserts compiled out ($833/833$ succeeded, $21$ min wall, $8$-way parallel).
Rows by hub degree: $d{=}7$: $40{,}257$; $d{=}8$: $8{,}070$; $d{=}9$: $93$;
$d{=}10$: $3$; $d{=}5$: $26$; $d{=}6$: $60$ (the $3{,}046$ low-degree wheels collapse
to only $86$ distinct rows). The $93$ $d{=}9$ and $3$ $d{=}10$ rows are exactly the
strict rows corresponding to the $51+2$ \texttt{degree\_range} failures of
\S\ref{sec:s1}.

\paragraph{Positive control.} The full-pool LP is feasible and $x_0$ satisfies every
one of the $3{,}511$ generated rows, with $0$ violations. The control was also run
without the $d\le 6$ rows to demonstrate that the trivially feasible regime exists
and is not what we report.

\paragraph{Why the headline is the audited certificate.} At $x_0$, the ring-$\le14$
system is violated in $9{,}647$ rows; but $x_0$ still violates $8{,}741$ rows even in
the fully conservative system, with $\mathrm{LB}$ up to $+3$. So at $x_0$ the
ring-$\le 14$ argument fails not merely as a verification artifact but genuinely, at
thousands of degree-$7$/$8$ hubs that really do retain positive charge. The
\texttt{-{}-exact-only} and \texttt{-{}-conservative} diagnostics lose the
contradiction, which is why the headline is carried by the audited $12$-row
irreducible subsystem rather than the raw $48{,}509$-row verdict.

\subsection{The $12$-row irreducible core and what it says}

The $19$-row Farkas support (Appendix~\ref{sec:certificate}) is not minimal. Greedy
deletion filtering shrinks it to a $12$-row \emph{irreducible infeasible subsystem}:
nine degree-$7$ hub rows and three degree-$5$ rows, irreducible in the strict sense
that deleting any one alone restores feasibility (verified for all $12$). Its
certificate has $y^\top b = -20$ with multipliers
$75, 75, 25, 75, 55, 60, 15, 55, 25, 75, 168, 219$. All four ring-$\le14$-only rows
survive into the core. This is the smallest human-inspectable object carrying the
whole result, and it is legible:

\begin{itemize}
  \item \textbf{No degree-$8$, $9$ or $10$ row appears.} The contradiction needs none
        of the $53$ \texttt{degree\_range} failures nor any of the $250$ degree-$8$
        failures observed empirically. This sharpens the empirical picture: the
        irreducible cause sits entirely at hub degree $7$, colliding with the
        degree-$5$ discharge.
  \item \textbf{Every one of the nine degree-$7$ hubs has a degree-$5$ spoke.} The
        whole contradiction lives in the degree-$5$/degree-$7$ interface.
  \item \textbf{Only $19$ of the $84$ rules carry a nonzero coefficient anywhere in
        the $12$ rows} (\texttt{rule001}, \texttt{rule002}$_{1,2}$,
        \texttt{rule003}$_{1,2}$, \texttt{rule004}$_{1,2}$, \texttt{rule008}$_{1,2}$,
        \texttt{rule009}$_{1,2}$, \texttt{rule010}$_{1,2}$, \texttt{rule011}$_{1,2}$,
        \texttt{rule023}$_{1,2}$, \texttt{rule032}$_{1,2}$).
        \textbf{$65$ of the $84$ shapes are irrelevant to the obstruction} ---
        re-tuning them cannot help.
  \item In the Farkas combination $78$ of the $84$ rule columns cancel exactly; only
        six carry slack. The near-total cancellation is why the contradiction is
        tight.
  \item Read as arithmetic: the degree-$5$ rows say a degree-$5$ vertex must give
        away all $10$ units of charge, forcing the amounts on
        \texttt{rule001}--\texttt{rule004} \emph{up}; the degree-$7$ rows cap what a
        degree-$7$ hub with degree-$5$ neighbours may receive back, pulling them
        \emph{down}. With the ring-$\le14$ pool no longer blocking four offending
        degree-$7$ cartwheels, the two demands cannot both be met.
\end{itemize}

\paragraph{The infeasibility is genuinely caused by the ring-$\le14$-only rows.}
Delete the four ring-$\le14$-only rows from the $19$-row support and the remaining
$15$ --- every one of which is also valid for the full pool --- are feasible, and
$x_0$ satisfies all $15$. The contradiction is not lurking in the full-pool
constraints; it is created by exactly the cartwheels the ring-$\le14$ pool fails to
block. Independently, those four rows come from wheels \texttt{d7\_3358},
\texttt{d7\_3372}, \texttt{d7\_3376} --- precisely three of the wheels whose job the
\emph{C++} aborted with \texttt{Assertion failed: (C == 0)}. The Python row semantics
and the C++ asserts agree wheel by wheel on the certificate's own support.

\subsection{Independent verification of the certificate}

\art{tools/verify\_farkas.py} is solver-free and uses \texttt{fractions.Fraction}
throughout --- no floating point anywhere on the checking path. Critically it does
not trust the certificate's stored coefficients: it re-derives every row from its
source (\texttt{.cartwheel} files re-parsed and re-run through \texttt{leaf\_row};
low-degree tags rebuilt with \texttt{generate\_cartwheel} and re-run through
\texttt{lowdeg\_row}) and fails loudly on any mismatch. It then checks
(a) $y\ge 0$, (b) $y\ne 0$, (c) $(y^\top A)_k \ge 0$ for all $84$ columns,
(d) $y^\top b < 0$, prints the contradiction chain
$0 \le (y^\top A)x = y^\top(Ax) \le y^\top b = -30 < 0$, and independently confirms
that $x_0$ violates $4$ of the $19$ support rows. Deliberate corruption of a single
coefficient is rejected with an exact claimed-versus-re-derived diff.

\subsection{The strong encoding, and why it is not used}
\label{sec:strongenc}

The alternative encoding --- \emph{every stage-1 wheel the pool does not block must
have $\texttt{upper\_bound\_of\_charge} < 0$} --- is what a naive reading suggests and
is what max-of-affine cutting-plane machinery is for. It is \emph{not} a correct
encoding of ``the argument closes'': survivors are permitted, and the published proof
at $x_0$ leaves $5{,}439 + 6{,}790 + 3{,}285 + 626 + 8 = 16{,}148$ stage-1 wheels
alive on purpose. Imposing it would reject the published proof itself.

We ran it anyway as a negative control, and it fails exactly as predicted: the
\emph{full-pool} run is \textsc{infeasible} at iteration $0$ with a verified $10$-row
certificate ($y^\top b = -175$; support: three degree-$5$, two degree-$6$ and five
degree-$7$ rows). So the strong condition cannot be met by any nonnegative amounts
even with the complete $8{,}200$-configuration pool --- it would ``disprove'' the
published theorem. That is the concrete demonstration that the strong encoding is the
wrong one. It is reported as a measurement and never as evidence for the headline.

\subsection{Scope}
\label{sec:lp-scope}

\begin{enumerate}
  \item \textbf{These $84$ shapes only.} Nothing rules out a different or larger set
        of rule shapes. Deciding that requires column generation over planar
        degree-interval patterns, whose pricing subproblem is a search over planar
        patterns rather than an LP (\S\ref{sec:open}).
  \item \textbf{This pool only.} $P_{14}$ is exactly the ring-$\le14$ part of the
        published pool. Many D-reducible configurations of ring $\le14$ are not in it.
  \item \textbf{The $d\le6$ rows are a reconstruction} (\S\ref{sec:lowdeg}); they are
        load-bearing, since without them the LP is feasible.
  \item \textbf{Unblocked $\ne$ realizable} (\S\ref{sec:notclaim}).
  \item \textbf{$x \ge 0$ is a hypothesis}, used in Lemma~\ref{lem:mono}(c) and in
        \eqref{eq:lowdeg}. Monotonicity was measured to hold on $2{,}145$
        parent/child pairs with $x\ge 0$ and to \emph{fail} on $8$ of them when
        negative amounts are allowed, so the hypothesis is exactly load-bearing, not
        decorative. In the other direction there is nothing to hedge: the paper fixes
        $r(R)\in\{1,2\}$ for all $84$ rules, and
        $\{1,2\}^{84}\subset\mathbb{Z}_{\ge0}^{84}\subset\mathbb{R}_{\ge0}^{84}$, so
        deciding infeasibility over the nonnegative reals is strictly stronger than
        the paper's own amount space requires.
  \item \textbf{Integrality is not needed}: all certificate rows are non-strict and
        the real variant is infeasible with the same certificate.
\end{enumerate}

\paragraph{Adversarial pass.} The formulation was handed to a reviewer prompted to
refute it, defaulting to ``the encoding is wrong'', across ten named attack lines.
Outcome: one real error found and fixed (the pool-independence claim in
Lemma~\ref{lem:mono}(a)); one genuine scope correction adopted (the code-necessary
versus math-necessary distinction, which became \S\ref{sec:necessity}); the rest
survived with evidence. A later pass against the paper's own LaTeX source found two
prose errors (both corrected: the paper does \emph{not} do the degree-$5$/$6$ case by
hand, and the e-print \emph{was} obtainable) and one genuine open gap (the coarse
cartwheel issue), which was closed exactly and exhaustively without changing the
certificate. Ablations: the degree-$6$ rows alone with all leaf rows are feasible;
the degree-$5$ rows alone are already infeasible. Six solver configurations all
report infeasible, and a uniform slack of $0.082$ on every right-hand side is the
minimum needed to restore feasibility.

%======================================================================
\section{$f(r)$: the minimum interior size of a D-reducible configuration}
\label{sec:fr}
%======================================================================

\subsection{Statement}

\begin{theorem}[Proven, exhaustive]
\label{thm:fr}
For each ring size $r \in \{8,9,10,11,12,13\}$, every D-reducible configuration of
ring size $r$ has at least $f(r)$ interior vertices, where
\[
\begin{array}{c|cccccc}
 r & 8 & 9 & 10 & 11 & 12 & 13\\\hline
 f(r) & 5 & 5 & 6 & 7 & 7 & 8
\end{array}
\]
and each bound is tight: a D-reducible configuration with exactly $f(r)$ interior
vertices exists in the published catalogs.
\end{theorem}

Birkhoff \cite{birkhoff1913} proved that a reducible configuration has ring size
$\ge 6$. To our knowledge Theorem~\ref{thm:fr} is the first published lower bound on
interior size as a function of ring size.

\subsection{The enumeration}
\label{sec:fr-enum}

For each $r$ and each $n$ from $r+1$ to $r+f(r)-1$ we generated \emph{every}
near-triangulation of the disk with outer face of size $r$ on $n$ vertices with
\texttt{plantri -P$r$ -c3 -a $n$}, filtered to the RSST configuration class
(nonempty connected interior of minimum degree $5$; RSST conditions 4--7 including
the ring-contact-arc condition 6), and ran the D-reducibility checker on each.
Table~\ref{tab:frsweep} reports the counts.

\begin{table}[t]
\centering\small
\caption{The exhaustive sweep. ``raw'' = near-triangulations from \texttt{plantri};
``cand'' = passing the interior filter; ``excl'' = additionally failing RSST
condition 6; ``valid'' = the true RSST configuration class. Every ``D-red'' entry is
$0$. Per-$(r,n)$ machine-readable counts in
\art{results/theorem/fr\_table/manifest\_r*.json}; per-configuration records
(adjacency, $a$, closure trace) in \art{configs\_r*\_n*.jsonl}.}
\label{tab:frsweep}
\begin{tabular}{rrrrrrr}
\toprule
$r$ & interior range & raw & cand & excl & \textbf{valid} & \textbf{D-red} \\
\midrule
8  & 1--4 & 436          & 18     & 6      & 12     & 0 \\
9  & 1--4 & 666          & 35     & 17     & 18     & 0 \\
10 & 1--5 & 13{,}338     & 221    & 127    & 94     & 0 \\
11 & 1--6 & 293{,}555    & 1{,}631& 1{,}028& 603    & 0 \\
12 & 1--6 & 509{,}288    & 4{,}092& 3{,}004& 1{,}088& 0 \\
13 & 1--7 & 12{,}389{,}976 & 36{,}392 & 27{,}888 & 8{,}504 & 0 \\
\midrule
\textbf{total} & & \textbf{13{,}207{,}259} & \textbf{42{,}389} & \textbf{32{,}070} & \textbf{10{,}319} & \textbf{0} \\
\bottomrule
\end{tabular}
\end{table}

Two remarks on the class definition, both of which cost real debugging.

\emph{RSST condition 6 is genuine structure and removes a lot.} After normalising
away the labelling convention (condition 4), a substantial fraction of accepted
near-triangulations still fail condition 6 --- some interior vertex meets the ring in
more than two contiguous arcs. \texttt{reduce.c}'s own \texttt{ReadConf} rejects such
graphs outright: they are not configurations in RSST's technical sense at all. The
exclusion rate rises from $33\%$ at $r=8$ to $77\%$ at $r=13$. Excluded graphs were
nevertheless retained in the output with a D-reducibility verdict computed anyway, so
that a D-reducible-but-not-RSST-legal graph would have been surfaced loudly. None
was: every one of the $32{,}070$ exclusions is non-D-reducible too, so the filter
never changed a verdict, only the population count.

\emph{Condition 4 is a labelling convention, not structure.} \texttt{ReadConf}
requires each ring vertex's neighbour list to start at a specific ring neighbour;
nothing our checker computes depends on it, since a cyclic list denotes the same
rotation system whatever entry it starts from. Naively serialising an RSST-relabelled
configuration made \texttt{build/reduce\_rsst} reject it with \texttt{ReadErr(4,\dots)}
even though the underlying embedded graph was fine. We added
\texttt{rotate\_ring\_starts}, a mathematical no-op on the rotation system.

The $r=13$ sweep is the expensive one: interior $6$ ($n=19$) and interior $7$
($n=20$) were run $8$-way sharded by residue class, $\approx 2{,}250$ s and
$\approx 13{,}300$ s per shard respectively, giving $1{,}445$ and $6{,}710$ valid
configurations and $0$ D-reducible.

\subsection{Triple verification}

\paragraph{(1) Our checker vs.\ the compiled 1995 oracle.} Up to $20$ RSST-legal
below-threshold configurations per ring were serialised and run through
\texttt{build/reduce\_rsst}, compiled from the pinned \texttt{reduce.c}: $70$ sampled,
\textbf{$0$ verdict mismatches}, $6$ oracle-inconclusive
(\art{results/theorem/fr\_table/crosscheck.json}).

\paragraph{(2) The six inconclusive cases, adjudicated by a third implementation.}
On $6$ of the $70$, the oracle recomputed a \emph{different} $|\Phi(K)|$ than our
checker supplied and aborted with \texttt{*** ERROR: DISCREPANCY IN NUMBER OF
EXTENDING COLOURINGS ***} before printing a verdict. We wrote a third, from-scratch
count of $|\Phi(K)|$ (\art{src/fourcolor/brute\_force.py}) sharing no code path with
the production checker --- different edge indexing, different triangle enumeration,
no gauge fixing, no greedy variable ordering. It agrees with our checker in all $6$
cases and disagrees with the oracle in all $6$:

\begin{center}\small
\begin{tabular}{lrrrrr}
\toprule
config & $r$ & $n$ & our $a$ & \texttt{reduce\_rsst} & brute force \\
\midrule
\texttt{fr-r8-n12-172}    & 8  & 12 & 72  & 48  & 72 \\
\texttt{fr-r9-n13-269}    & 9  & 13 & 147 & 94  & 147 \\
\texttt{fr-r10-n15-4669}  & 10 & 15 & 367 & 230 & 367 \\
\texttt{fr-r10-n15-4678}  & 10 & 15 & 365 & 243 & 365 \\
\texttt{fr-r10-n15-4613}  & 10 & 15 & 408 & 250 & 408 \\
\texttt{fr-r11-n17-89446} & 11 & 17 & 995 & 532 & 995 \\
\bottomrule
\end{tabular}
\end{center}

\emph{Provenance note.} Our own and the brute-force counts, and the
\textsc{oracle\_inconclusive} status of exactly these six, are machine-readable in
\art{results/theorem/fr\_table/crosscheck.json} (which records
\texttt{oracle\_returncode: 31} for each). The oracle's \emph{recomputed} counts in
the middle column were transcribed from its \texttt{stdout} into
\art{results/theorem/fr\_table/THEOREM.md} \S4.3 and are not retained in the JSON;
they are the one set of numbers in this paper that is not re-derivable from a
machine-readable artifact without re-running the oracle.

Our checker agrees with the oracle's header $a$ on all $3{,}455$ hand-curated
catalog configurations with zero mismatches (\S\ref{sec:verify}), so this is not a
general defect. The evidence points to \texttt{reduce.c}'s \texttt{strip()}
edge-numbering heuristic --- a hand-optimised greedy ordering routine, documented in
its own comments as chosen to shrink the search tree --- being the outlier on a class
of inputs (small, exhaustively generated, ``generic'' near-triangulations) its
original validation against a hand-curated catalog never exercised. None of the six
configurations' D-reducibility verdict is in question: all are \texttt{False}, and
the oracle never disputed the verdict, only the intermediate count, before erroring
out. We report it because an honest cross-verification protocol surfaces every
discrepancy it finds.

\paragraph{(3) Tightness witnesses.} For each $r$, a catalog D-reducible
configuration with exactly $f(r)$ interior vertices, verified both by our checker and
by the oracle (\art{results/theorem/fr\_table/tightness\_witnesses.json}):

\begin{center}\small
\begin{tabular}{rlrrrl}
\toprule
$r$ & witness & $n$ & interior & our $a$ & oracle \\
\midrule
8  & \texttt{D0003.conf} (nl4ct D) & 13 & 5 & 100   & \texttt{*** D-reducible ***} \\
9  & \texttt{D0006.conf} (nl4ct D) & 14 & 5 & 211   & \texttt{*** D-reducible ***} \\
10 & \texttt{D0015.conf} (nl4ct D) & 16 & 6 & 536   & \texttt{*** D-reducible ***} \\
11 & \texttt{D0024.conf} (nl4ct D) & 18 & 7 & 1{,}363 & \texttt{*** D-reducible ***} \\
12 & \texttt{D0143.conf} (nl4ct D) & 19 & 7 & 2{,}863 & \texttt{*** D-reducible ***} \\
13 & \texttt{126.504366} (RSST 633) & 21 & 8 & 6{,}954 & --- \\
\bottomrule
\end{tabular}
\end{center}

Across the whole labelled corpus, the minimum interior size of a D-reducible
configuration is $4,4,5,5,6,7,7,8,9,10,10$ at $r=6,\dots,16$; the entries at
$r=8,\dots,13$ are exactly $f(r)$, and the entries at $r=14,15,16$ are upper bounds
on $f$ there for which we have no exhaustive lower-bound sweep.

\subsection{The C-versus-D nuance}
\label{sec:cvsd}

Theorem~\ref{thm:fr} is scoped to D-reducibility (no contract). C-reducibility is
strictly weaker, and the RSST $633$-configuration catalog itself contains
configurations at ring sizes $8$ and $11$ with interior sizes \emph{strictly below}
$f(r)$ that are C-reducible but not D-reducible:

\begin{center}\small
\begin{tabular}{rlrrl}
\toprule
ring & catalog ident & $n$ & interior & contract \\
\midrule
8  & \texttt{2.126}    & 12 & 4 & $\{(1,9),(3,9),(5,11),(7,11)\}$ \\
11 & \texttt{126.7566} & 17 & 6 & $\{(15,12)\}$ \\
11 & \texttt{128.136}  & 17 & 6 & $\{(15,12),(1,12),(3,12),(10,15)\}$ \\
11 & \texttt{7562.126} & 17 & 6 & $\{(15,12)\}$ \\
\bottomrule
\end{tabular}
\end{center}

So the corresponding minimum over all reducible (C-or-D) configurations is, at least
at rings $8$ and $11$, strictly smaller, and the D-qualifier in
Theorem~\ref{thm:fr} is not decoration. (No ring-$9$, $10$ or $12$ catalog member
exists at interior $f(r)-1$ in either catalog; that is a fact about which
configurations the catalogs' authors happened to need, not evidence that C-reducible
near-misses do not exist at other rings.)

\subsection{A first out-of-sample test of the mass law}
\label{sec:gap}

The $f(r)$ sweep produces, as a by-product, an exhaustive population of
\emph{below-threshold non-reducible} configurations with their exact $a$ values.
Comparing the maximum $a$ over that population with the minimum $a$ over the
D-reducible population at the same ring gives a clean separation at every measurable
ring:

\begin{center}\small
\begin{tabular}{rrrr}
\toprule
$r$ & $\max a$ (below threshold, non-red.) & $\min a$ (D-reducible) & gap \\
\midrule
8  & 82      & 94      & 12 \\
9  & 172     & 211     & 39 \\
10 & 438     & 496     & 58 \\
11 & 1{,}128 & 1{,}252 & 124 \\
12 & 2{,}339 & 2{,}863 & 524 \\
13 & 6{,}003 & 6{,}954 & 951 \\
\bottomrule
\end{tabular}
\end{center}

The $r=13$ row is genuinely out of sample: the prediction (a gap) was made after the
$r\le12$ data and before the $r=13$ enumeration finished. This is the empirical
content of the threshold half of the coloring-mass law (\S\ref{sec:mass}).

%======================================================================
\section{Independent verification of the 4CT corpus}
\label{sec:verify}
%======================================================================

Nothing in \S\ref{sec:lp}--\S\ref{sec:fr} means anything unless our reducibility
checker is right. This section reports what we did to establish that, and what we
found while doing it. The checker \texttt{fourcolor.reduce.check} was written from
the mathematics of \cite{rsst-reduce} (the \texttt{reduce.tex} write-up, in
particular its Theorem 3.2 machinery), \emph{not} ported from \texttt{reduce.c}; so
agreement with the C programs is genuine differential evidence rather than a
tautology.

\subsection{The three catalogs}

\begin{table}[h]
\centering\small
\caption{Independent re-verification of the published catalogs. ``Agree'' means the
computed $|\Phi(K)|$ and $|{\rm consistent}|$ both match the published header fields
\emph{and} the configuration is confirmed reducible.}
\label{tab:catalogs}
\begin{tabular}{llrrl}
\toprule
catalog & source & configs & agree & artifact \\
\midrule
RSST $633$ & \cite{rsst-reduce} \texttt{unavoidable.conf} & 633 & \textbf{633} &
  \art{results/differential-unavoidable-r14/report.jsonl} \\
Steinberger $\mathcal{U}_{2822}$ & \cite{steinberger2010} \texttt{U\_2822.conf} & 2{,}822 & \textbf{2{,}822} &
  \art{results/differential-U\_2822-r16/report.jsonl} \\
2026 pool (pruned) & \cite{nl4ct2026} \texttt{D/*.conf} & 7{,}770 & \textbf{7{,}697} &
  \art{results/p3/pool-verify/SUMMARY.json} \\
\bottomrule
\end{tabular}
\end{table}

Three points of precision that we have seen stated loosely elsewhere, including in
our own earlier notes.

\paragraph{The RSST $633$ are not all D-reducible.} Exactly $249$ are D-reducible;
the other $384$ are C-reducible via a nonempty published contract and are
\emph{not} D-reducible. The compiled oracle confirms the same split
($249$ \texttt{*** D-reducible ***}, $384$ \texttt{*** Contract confirmed ***}
alongside $384$ \texttt{*** Not D-reducible ***}). What is verified for all $633$ is
the header fields and reducibility in the C-or-D sense. Ring sizes run $6$--$14$,
so no configuration was excluded by our ring bound.

\paragraph{Steinberger's set is genuinely D-only.} All $2{,}822$ are D-reducible,
all with header $b=0$ (no contract), rings $6$--$16$; $2{,}822/2{,}822$ agree with
the published headers, $0$ disagreements, $13{,}103$ s of checker CPU. This matters
because \S\ref{sec:lowdeg}'s citation chain rests on Steinberger's degree-$5$/$6$
computation being against a correct set.

\paragraph{The 2026 pool is verified for $7{,}697$ of $8{,}200$.} The verified run
used the pruned pool of $7{,}770$ (see below); of those, $7{,}697$ have ring $\le16$
and were verified D-reducible with $0$ anomalies in $12.3$ CPU-hours, and $73$ at
rings $17$--$18$ are outside our engines' ring budget and \textbf{remain unverified
by us}. Together with the $430$ removed configurations, the honest accounting is
$8{,}200 = 7{,}697\ \text{verified} + 430\ \text{removed} + 73\ \text{unverified}$.
Prior to this, the 2026 pool's reducibility had never been independently verified:
the authors' own verification was external to the released code.

\subsection{A legacy oracle that aborts, and what replaced it}

We compiled two C oracles: \texttt{build/reduce\_rsst} from \cite{rsst-reduce}'s
\texttt{reduce.c}, and \texttt{build/reduce\_stein} from the Steinberger-extended
version (needed because the unmodified 1995 code has \texttt{MAXRING} $=14$).

\texttt{reduce\_rsst} completes the $633$ cleanly
(\texttt{Reducibility of 633 configurations verified}). \texttt{reduce\_stein}
does \emph{not} complete Steinberger's $2{,}822$: it processes $783$ of them
(of which $782$ print \texttt{*** D-reducible ***}) and then aborts with
\texttt{SIGABRT} (exit $134$), truncating mid-number. The crashing input is
preserved as \art{build/crash783.conf}: a ring-$16$ configuration with
$n=28$ and $a=139{,}480$. Our checker completes it in $227.7$ s with
$a = 139{,}480$, matching the published header. So the full-set verification of
$\mathcal{U}_{2822}$ rests on our checker against the published headers, not on that
binary; the binary corroborates the first $783$ and then dies.

We also encountered the milder failure of \S\ref{sec:fr}: on six exhaustively
generated near-triangulations, \texttt{reduce\_rsst} recomputes $|\Phi(K)|$
incorrectly and self-aborts, and a third implementation adjudicates in our favour in
all six. The two findings together are the practical reason the programme maintains
\emph{three} independent counts of $|\Phi(K)|$ rather than two.

\subsection{Reproducing the 2026 pipeline}

We reproduced the released C++ pipeline \cite{nl4ct2026} end to end on the full
$8{,}200$-configuration pool ($19{,}754$ configurations after symmetry expansion) and
the $84$ rule shapes, in a Debug build with all assertions live. Total wall time
$4$\,h\,$38$\,m. Every count matches the paper:

\begin{center}\small
\begin{tabular}{llr}
\toprule
stage & result & wall \\
\midrule
\texttt{combine\_rules}   & $1{,}832$ combined rules, $671$ non-blocked & $\approx3$ s \\
\texttt{enum\_wheels} $d{=}7..11$ & $5{,}439 / 6{,}790 / 3{,}285 / 626 / 8$ bad wheels & $1$\,h\,$45$\,m \\
\texttt{enum\_cartwheels} & $10{,}094$ bad cartwheels ($9{,}366$ at $d{=}7$, $728$ at $d{=}8$; zero at $d{=}9,10,11$) & $2$\,h\,$09$\,m \\
three gluing lemmas & all passed ($5{,}365$ / $1{,}618$ / $298$ cartwheels) & $43$ m \\
\bottomrule
\end{tabular}
\end{center}

Separately we re-implemented the wheel-level stage in Python
(\art{src/fourcolor/nl4ct.py}) and ran it differentially against the C++ on
\textbf{all $5{,}692{,}937$ wheel candidates} at hub degrees $7$--$11$
($11{,}165 + 48{,}915 + 217{,}045 + 976{,}887 + 4{,}438{,}925$): every degree reports
\texttt{exact\_match: true} with $0$ false positives and $0$ false negatives against
the C++ ground truth. That validation is what licenses the symbolic charge form
\eqref{eq:leafrow} used in \S\ref{sec:lp}.

We also replayed the RSST 1995 unavoidability argument: \texttt{discharge.c} from
\cite{rsst-discharge} against the five presentation files \texttt{present7}
through \texttt{present11}, all five verified against the $633$-configuration set
(\art{results/discharge-rsst/replay.txt}).

\subsection{Coverage measurement and a pool that is $430$ configurations too large}

Measuring which pool configurations the pipeline actually uses gives a sharply
two-level picture. At the \emph{wheel} level, only $85$ of the $8{,}200$
configurations ever act as a blocker anywhere across $d=7..11$; the whole wheel-level
covering structure has at most one concretization per wheel and at most $5$ blockers
per concretization. At the \emph{cartwheel} (refinement) level, $7{,}422$ distinct
configurations appear as blockers over $2{,}253{,}185$ logged blocking events ---
$90.5\%$ of the pool, and that is a \emph{lower} bound, because the instrumentation
records only the first matching configuration per event. So the intuition that a
small core carries the proof is right at stage 1 and wrong at stage 2.

We then asked which configurations are used by \emph{no} stage. A first attempt
tracked three roles (combined-rule blocking, wheel blocking, first-match cartwheel
blocking) and found $673$ unused; removing them made the pipeline fail --- all three
gluing checks aborted, and the bad-cartwheel count collapsed from $10{,}094$ to $0$.
That failure is preserved (\art{results/p3/runs/batch1/}) and it is what discovered
a \emph{fourth} blocking role: the gluing lemmas also consult the pool ($2{,}199$
distinct glue-blockers). Tracking all four roles gives \textbf{$430$ configurations
with zero usage everywhere}. Removing exactly those and re-running the entire
pipeline gives \textsc{pass} with counts \emph{identical to the full-pool proof in
every field}: $1{,}832/671$ combined rules, $5{,}439/6{,}790/3{,}285/626/8$ wheels,
$10{,}094$ bad cartwheels, three gluing lemmas green
(\art{results/p3/runs/batch1b/VERDICT.txt}, \art{counts.json}).

\begin{proposition}[Certified]
The 2026 proof goes through verbatim with a pool of $7{,}770$ configurations rather
than $8{,}200$, a reduction of $5.2\%$.
\end{proposition}

This is a floor we reached, not a claimed minimum: no further prune batch was run,
and the honest reading of the two-level coverage picture above is that further
reduction is bounded by the cartwheel-level usage, which is nearly the whole pool.

\begin{remark}[What was designed and not built]
Our original plan (\art{01-PROBLEM-STATEMENT.md}) was a set-cover MILP minimising
$|U|$, with an LP-dual lower bound. \textbf{No MILP was ever built or solved}; the
design was replaced by batched prune-and-verify once the coverage measurement showed
that the first-match cartwheel attribution makes a one-shot set cover unsound. We
report this because the problem statement is in the repository and a reader might
otherwise look for a solver artifact that does not exist. Theorem~\ref{thm:lp} is
the useful residue of that line: an LP result, but an impossibility rather than an
optimisation.
\end{remark}

%======================================================================
\section{The coloring-mass law}
\label{sec:mass}
%======================================================================

\subsection{Setting}

For a configuration $K$ with free completion $S$, ring size $r$, interior size $k$,
$n=r+k$, write $a = a(K) = |\Phi(K)|$ and $W(r) := (2^r+2)/6$. The organising
conjecture of this section has two halves.

\begin{conjecture}[Coloring-Mass Law]
\label{conj:mass}
\emph{(T) Threshold.} There is an increasing $A(r)$ such that every D-reducible
configuration of ring size $r$ has $a \ge A(r)$; empirically
$A(r) \approx 2.4^{\,r}/12.79$.
\emph{(C) Cap.} $a \le B(r,k)$ with $B$ increasing in $k$; conjecturally the
\emph{sharp cap} $B(r,k) = W(r)\,(4/3)^{k-1}$.
\end{conjecture}

If both held, one would obtain analytic lower bounds on $f(r)$ for all $r$: the
finite table of Theorem~\ref{thm:fr} would become an infinite structural law
--- \emph{reducibility is a coloring-mass threshold phenomenon, and interior vertices
supply the mass}.

The Bridge Lemma below is the bridge from colourings to the chromatic polynomial and
underlies everything that follows.

\begin{lemma}[Bridge Lemma; proven]
\label{lem:bridge}
For every near-triangulation $S$ of the disk, the number of tri-colourings of $S$ is
$P(S,4)/4$; the $S_3$ action permuting the three edge colours is free on
tri-colourings, so the number of tri-colouring classes is $P(S,4)/24$; and
consequently $a \le P(S,4)/24$.
\end{lemma}

\begin{proof}[Proof sketch]
Identify the three edge colours with the nonzero elements of the Klein group
$V=\mathbb{Z}_2^2$ and label each edge $uv$ of a proper $4$-colouring $c$ by
$t(uv)=c(u)+c(v)$. Properness makes $t$ a tri-colouring, and translation by $V$ gives
a $4$-to-$1$ map. Conversely a tri-colouring is a $V$-valued $1$-cochain summing to
$0$ around every internal face; since $S$ is a triangulated disk the internal faces
generate the cycle space, so $t$ is exact and has exactly $4$ potentials. For
freeness: a transposition fixing $t$ would force all edges to carry the one fixed
colour, impossible in a triangle; a $3$-cycle fixes no colour. Restriction to the ring
sends each class of tri-colourings to one class of ring colourings, and $\Phi(K)$ is
the image.
\end{proof}

\subsection{The wheel: exact closed form}

\begin{theorem}[Proven, with equality case]
\label{thm:wheel}
If $k=1$ then $S$ is the wheel $W_r$ and
\[
   a \;=\; \frac{2^r + 2(-1)^r}{6} \;=\; \frac{P(W_r,4)}{24}.
\]
In particular the Bridge Lemma is \emph{tight} at every wheel; $a = W(r)$ for even
$r$ and $a = \lfloor W(r)\rfloor$ for odd $r$.
\end{theorem}

\begin{proof}
$k=1$ forces the internal faces to be $h x_i x_{i+1}$, so $S=W_r$, and interior
degree $\ge5$ forces $r\ge5$. Writing $e_i$ for the rim edge colour and $s_i$ for the
spoke colour, the face condition says $\{s_i,e_i,s_{i+1}\}$ is all three colours, so
$s_i\ne s_{i+1}$ and $e_i$ is determined; hence
$\#\mathrm{TriCol}(W_r) = P(C_r,3) = 2^r+2(-1)^r$.
Let $\rho$ be the restriction to the rim. Given $(e_i)$ and $s_1$, the face condition
forces $s_{i+1} = \tau_{e_i}(s_i)$ where $\tau_c$ transposes the two colours $\ne c$;
so the spoke vector is determined by $s_1$, which must avoid $e_r$ and $e_1$. If
$e_i \ne e_{i+1}$ for some $i$ then $s_{i+1}$ avoids two distinct colours and the
fibre has size $\le1$. If $(e_i)$ is monochromatic the spokes alternate between the
two colours $\ne c$, which closes up iff $r$ is even, giving a fibre of size exactly
$2$ (and empty for odd $r$). An element of $S_3$ fixes $(e_i)$ iff it fixes every
$e_i$; a $3$-cycle fixes no colour and a transposition fixes one, so the rim
colourings with nontrivial stabiliser are exactly the three monochromatic ones.
For odd $r$ no monochromatic rim colouring is in the image, every fibre is a
singleton, $S_3$ acts freely, and $a = |\mathrm{TriCol}|/6=(2^r-2)/6$. For even $r$
the six tri-colourings above the three monochromatic rim colourings form a single
free $S_3$-orbit --- one class --- while the three rim colourings likewise form one
element of $\Phi$, so $a = (|\mathrm{TriCol}|-6)/6 + 1 = (2^r+2)/6$.
\end{proof}

The wheel value is $P(C_r,3)/6$, the number of proper $3$-colourings of a cycle
modulo colour symmetry. Everything in the sharp cap's constant is that one classical
number. The theorem was adversarially re-verified four independent ways for
$r=5..14$ --- our production checker, the brute-force tie-breaker, a from-scratch
enumeration of spoke sequences with explicit $S_3$ canonicalisation, and the compiled
1995 RSST oracle (which covers $r \le 12$ and hard-errors at $r\ge13$ on a
compile-time array cap, a build limit rather than a disagreement).

\subsection{The star-first cap}

\begin{theorem}[Proven]
\label{thm:star}
For every configuration $K$ with $k\ge1$,
\[
   a \;\le\; \min_{h \text{ interior}} \frac{(2^{d}+2(-1)^{d})\,2^{\,n-d-1}}{6},
   \qquad d = \deg h,
\]
and since every interior degree is $\ge 5$, uniformly
\[
   \boxed{\;a \;\le\; 11\cdot 2^{\,r+k-7} \;=\; \tfrac{11}{16}\,2^{\,r+k-3}.\;}
\]
The bound is exactly attained at every wheel.
\end{theorem}

\begin{proof}[Proof sketch]
Let $h$ be interior of degree $d$ and $L=N(h)$ its link, a cycle (not necessarily
induced --- a chord only removes colourings, so $P(C_d,3)$ remains valid). Order the
vertices $h, x_1,\dots,x_d, w_1,\dots,w_{n-d-1}$ so that every $w_j$ has two
predecessors adjacent to each other: grow greedily from the closed star of $h$ by
attaching any triangle sharing an edge with the built region. Such a triangle
introduces at most one new vertex, adjacent to both ends of the shared edge. The
greedy exhausts the rest because the dual graph of $S$ is connected and, for any
triangle $T$ not incident with $h$, the last triangle on a dual path from $T$ to a
star triangle shares a \emph{link} edge $x_i x_{i+1}$ (it cannot share $hx_i$, which
lies only in star triangles). Counting along the order: $h$ has $4$ choices, $L$ has
at most $P(C_d,3)=2^d+2(-1)^d$, and each $w_j$ has at most $2$. Divide by $24$ using
Lemma~\ref{lem:bridge}. For the uniform form, $2^{6-d}(2^d+2(-1)^d)/6$ equals
$10, 11, 10.5, 10.75, 10.625, \dots \to 32/3$ for $d = 5,6,7,8,9,\dots$, with maximum
$11$ at $d=6$.
\end{proof}

Note the counter-intuitive shape: \emph{odd} interior degrees give the better bound,
since $P(C_d,3) = 2^d-2$ for odd $d$. Consequently the worst case for the star-first
argument is a configuration whose interior degrees are all even and $\ge6$; and if
$S$ has an interior vertex of degree exactly $5$ (which $58{,}977$ of our $59{,}142$
corpus configurations do) then $a \le 5\cdot 2^{r+k-6}$. The $165$ exceptions are
all-degree-$\ge6$ interiors including the six wheels, which is exactly why the
unconditional constant must be $11$ rather than $10$.

Theorem~\ref{thm:star} improves the constant of the previously proven weak cap
$a \le 2^{r+k-3}$ by a factor $11/16$ and is exactly tight at $k=1$; but its
$k$-dependence is still $2^k$, so it moves the offset of the implied $f(r)$ bound and
not its slope: $f(r) \ge \lceil 0.26303\,r - 0.1367\rceil$ instead of
$\lceil 0.26303\,r - 0.6768\rceil$. That is $+1$ interior vertex at $6$ of $10$
rings, real, small, and free.

\subsection{$\gamma = 3/2$ at $k = 2$, and why the induction stops}

Order the interior graph by BFS from a vertex $h_1$ of maximum degree $d_1$. Colour
$h_1$ and its link; for $j\ge2$, split $\mathrm{link}(h_j)$ into maximal arcs of
not-yet-coloured vertices, each an at-most-$3$-colour path with both ends constrained,
so an arc of length $t$ admits at most $M(t) = 2,3,6,11,22,43,\dots$ colourings.
Writing $g(t)=M(t)/2^t \in \{1, 3/4, 3/4, 11/16, \dots\} \to 2/3$ gives the
per-configuration certificate
\begin{equation}
\label{eq:bfscert}
  a \;\le\; \bigl(1 + 2(-1)^{d_1}2^{-d_1}\bigr)\cdot\tfrac{2}{3}\cdot 2^{\,n-3}
      \cdot \textstyle\prod_{\text{arcs}} g(t),
\end{equation}
which reduces exactly to Theorem~\ref{thm:wheel} at $k=1$.

\begin{theorem}[Proven, $k=2$ only]
\label{thm:gamma32}
For $k=2$,
$\;a \le \bigl(1+2(-1)^{d_1}2^{-d_1}\bigr)\,\frac{2^r}{6}\cdot\frac{3}{2}
      \le \frac{33}{32}\cdot\frac{2^r}{6}\cdot\frac{3}{2}$.
\end{theorem}

\begin{proof}
Every edge of $S$ joining two interior vertices lies in exactly two triangles, hence
has at least two common neighbours; a third common neighbour $z$ would make
$h_1h_2z$ a separating triangle, which at $k=2$ is impossible since the only vertices
are $h_1,h_2$ and the ring, and the ring lies outside. So the coloured part of
$\mathrm{link}(h_2)$ at step $2$ is exactly $\{h_1,x,y\}$, the complementary arc is a
single arc of $\deg(h_2)-3 \ge 2$ vertices, and $\prod g = g(\deg(h_2)-3) \le 3/4$.
Substituting in \eqref{eq:bfscert} with $2^2\cdot(3/4)=3$ gives the claim.
\end{proof}

The ``exactly two common neighbours'' fact was machine-checked on all $610{,}846$
interior--interior edges of the corpus with zero exceptions, but it is only
\emph{proved} at $k=2$.

\paragraph{Why the induction stops, precisely.} The natural generalisation ---
hypothesis (H2), that every $h_j$ with $j\ge2$ contributes at least one arc of length
$\ge2$ --- would give $\gamma=3/2$ for all $k$. \textbf{(H2) is false on the corpus.}
Of the $372{,}008$ BFS steps with $j\ge2$, $20{,}379$ ($5.48\%$) contribute no arc at
all and $119{,}705$ ($32.18\%$) contribute only length-$1$ arcs, which gain nothing
since $g(1)=1$; the length-$1$-only fraction climbs from $\approx3\%$ at $k=3$ to
$\approx60\%$ at $k=16$. The realised effective base
$\gamma_{\text{prov}} = 2(\prod g)^{1/(k-1)}$ has maximum
$1.5000, 1.7321, 1.8171, 1.8612, 1.8882, 1.9064, 1.9195$ at $k=2,\dots,8$ and
$1.937082$ overall. So this scheme yields \emph{no uniform $\gamma<2$ beyond $k=2$}.
The value $1.7321 = 2\sqrt{3/4}$ at $k=3$ is the signature of ``one gaining step out
of two''.

Moreover $g(t)\ge 2/3$ for every $t$, so \emph{no argument of this shape can ever
give $\gamma < 4/3$}: the empirical near-tightness of $4/3$ is not an accident, it is
the scheme's infimum, attained only for arbitrarily long arcs.

\subsection{The obstruction theorem}

\begin{theorem}[Proven; explicit $10$-vertex witness]
\label{thm:obstruction}
Let $S_0$ be the free completion of the corpus configuration
\texttt{gen-r8-n10-2} ($r=8$, $k=2$, both interior vertices of degree $6$). Then
exactly
\[
   P(S_0,4) = 1464, \qquad \frac{P(S_0,4)}{24} = 61, \qquad a(S_0) = 55,
   \qquad W(8)\cdot\tfrac43 = \tfrac{172}{3} = 57.33\ldots
\]
Hence $P(S_0,4)/24 > W(8)\cdot(4/3) > a(S_0)$, so \textbf{no upper bound on $a$ that
factors through the Bridge Lemma --- i.e.\ any proof of the form
$a \le P(S,4)/24 \le B(r,k)$ --- can prove the sharp cap}, and more precisely no such
proof can establish $a \le W(r)\gamma^{k-1}$ for any
$\gamma < 61/43 = 1.4186\ldots$
\end{theorem}

Over the whole corpus the smallest $\gamma$ for which the $P/24$ form survives is
$\max\gamma_P = 1.4186047$ at exactly this record, against $\max\gamma_a = 1.3224310$
at \texttt{fr-r13-n15-4}; and the $P$-form of the sharp cap is outright \emph{killed},
failing on $1{,}960$ of $59{,}142$ records ($3.3\%$), worst ratio $1.3707$.

The Bridge-Lemma fibre $P(S,4)/(24a)$ --- two colour-inequivalent tri-colourings
sharing a ring restriction --- is exactly $1$ at $k=1$ and grows like $1.05$--$1.08^k$.
It is that factor the sharp cap silently uses. A proof of the sharp cap must
therefore contain a mechanism that sees \emph{ring restrictions}, not just
tri-colourings.

\paragraph{The per-step barrier.} There is a second, independent reason $4/3$ is hard.
For any order in which each vertex has $b_i \ge 2$ earlier neighbours, the edge count
$m = 3n-r-3$ gives the exact identity $\sum_{i\ge4}(b_i-2) = k$: the ``excess
adjacency'' budget is exactly one closing per interior vertex. A closing step with
earlier neighbours forming a path $u\!-\!x\!-\!w$ has $2$ colours available if
$\varphi(u)=\varphi(w)$ and $1$ otherwise, i.e.\ an average factor
$1 + \Pr[\varphi(u)=\varphi(w)]$. So the cap base is exactly $\gamma = 1+\Pr$, and
$\gamma=4/3$ means $\Pr \le 1/3$, i.e.\ $u$ and $w$ \emph{exactly uncorrelated} given
their common neighbour. That is a sharp mean-field statement at a critical point: the
$q=4$ antiferromagnetic Potts model on the triangular lattice is critical
\cite{moore2000}, correlations do not decay exponentially, and the standard
correlation-decay techniques \cite{salas1997} need $q\ge11$ there. Indeed
$4/3 = (q-2)^2/(q-1)$ at $q=4$ is Shrock--Tsai's \emph{lower} bound for the
triangular-lattice entropy \cite{shrock1997} --- a floor, not a ceiling.

We measured the distribution directly. Over a stratified sample of $1{,}315$
configurations covering every occupied $(r,k)$ cell, the exact step factors
$f_i = P(G_i,4)/P(G_{i-1},4)$ along the star-first order are:

\begin{center}\small
\begin{tabular}{rrrrr}
\toprule
$b_i$ & mean $f_i$ & min & max & steps \\
\midrule
2   & 2.0000 & 2.0000 & 2.0000 & 14{,}886 \\
3   & 1.3155 & 1.1754 & \textbf{1.6479} & 6{,}322 \\
4   & 0.8583 & 0.7273 & 1.0933 & 1{,}165 \\
$\ge5$ & 0.5138 & 0.2148 & 0.6576 & 565 \\
\bottomrule
\end{tabular}
\end{center}

The $b_i=3$ row is exactly $1+\Pr[\varphi(u)=\varphi(w)]$. Its mean $1.3155$ sits
just below the mean-field $4/3$ --- which is why the sharp cap fits so well on average
--- but its \emph{maximum} is $1.6479$, i.e.\ $\Pr = 0.648 > 1/2$. So no per-closing-step
bound can give $\gamma=3/2$, let alone $4/3$: individual closings are demonstrably
worse than either target, and any proof must average over them.

\subsection{The sharp cap survives where the Bridge Lemma does not}

Baxter's ground-state entropy of the $4$-state antiferromagnetic Potts model on the
triangular lattice is $\approx 1.4610$ per site \cite{baxter1987}, greater than
$4/3$. A disk cut from the triangular lattice has all interior degrees $6 \ge 5$ and
an induced boundary cycle, so it lies in our class --- and it predicts that the
$P$-form of the cap must fail badly in the bulk. It does: on exact $P(S,4)$
computations for hexagonal disks up to $R=7$ ($n=169$) and corner-truncated
triangular disks up to $m=16$, the ratio $(P/24)/\text{cap}$ reaches
$1.1\times10^4$, with fitted per-site entropy $1.541$/$1.563$.

But $a=|\Phi|$ itself, computed directly on the same disks by enumerating proper
$4$-colourings, taking Klein edge labels, restricting to the boundary cycle and
canonicalising under $S_3$ (self-validated $12/12$ against the repository's stored
$a$), \textbf{never violates the sharp cap, and the margin improves as the disk
grows}:

\begin{center}\small
\begin{tabular}{lrrrrrr}
\toprule
disk & $r$ & $k$ & $a$ & cap & $a/\text{cap}$ & $\gamma_a$ \\
\midrule
HEX $R{=}1$ ($=W_6$) & 6 & 1 & 11 & 11.0 & \textbf{1.0000} & --- \\
TRI $m{=}4$ & 9 & 3 & 140 & 152.3 & 0.9193 & 1.278 \\
TRI $m{=}5$ & 12 & 6 & 2{,}263 & 2{,}878.2 & 0.7863 & 1.271 \\
HEX $R{=}2$ & 12 & 7 & 2{,}756 & 3{,}837.5 & 0.7182 & 1.262 \\
TRI $m{=}6$ & 15 & 10 & 41{,}971 & 72{,}740.1 & 0.5770 & 1.255 \\
\textbf{TRI $m{=}7$} & \textbf{18} & \textbf{15} & \textbf{853{,}124} & \textbf{2{,}452{,}078} & \textbf{0.3479} & \textbf{1.2365} \\
\bottomrule
\end{tabular}
\end{center}

The Bridge-Lemma fibre on these disks is $1.21, 1.65, 1.93, 2.81, 6.30$ at
$k=3,6,7,10,15$: it grows fast enough to absorb the entire bulk excess and then some.
These disks were constructed adversarially to break the sharp cap and did not. What
is dead is the route, not the statement.

\subsection{The threshold half, empirically}

Mining the corpus for the exact frontier gives, per ring, the minimum $a$ over
D-reducible configurations:

\begin{center}\small
\begin{tabular}{rrrrl}
\toprule
$r$ & \#D-red & $\min a$ & $k^\ast$ & witness \\
\midrule
6  & 28      & 16       & 4  & \texttt{0.7322} (RSST 633) \\
7  & 37      & 39       & 4  & \texttt{2.122} (RSST 633) \\
8  & 1{,}435 & 94       & 5  & \texttt{gen-r8-n13-1494} (generated) \\
9  & 2{,}168 & 211      & 5  & \texttt{2.7566} (RSST 633) \\
10 & 1{,}891 & 496      & 6  & \texttt{122.140} (RSST 633) \\
11 & 209     & 1{,}252  & 7  & \texttt{122.1028} (RSST 633) \\
12 & 767     & 2{,}863  & 7  & \texttt{126.453966} (RSST 633) \\
13 & 1{,}961 & 6{,}954  & 8  & \texttt{126.504366} (RSST 633) \\
14 & 2{,}654 & 17{,}440 & 9  & \texttt{2039} (Steinberger) \\
15 & 1{,}562 & 42{,}957 & 10 & \texttt{2420} (Steinberger) \\
16 & 457     & 100{,}454& 10 & \texttt{D3848} (nl4ct pool) \\
\bottomrule
\end{tabular}
\end{center}

Two structural observations. First, \textbf{at every one of the eleven rings the
minimum-$a$ witness sits at the minimum possible interior count} --- on rings
$8$--$13$ that is exactly the $f(r)$ table. The coloring-mass frontier and the
interior-count frontier are attained by the \emph{same} configurations; the mass law
is not an independent phenomenon layered on top of Theorem~\ref{thm:fr}. Second, the
residual $\min a(r)/2.4^{\,r}$ is flat: $0.07823$ to $0.08540$ over eleven rings, a
spread of $1.0916$ across a range in which $2.4^r$ itself grows by a factor $6{,}340$.
Sweeping the base in steps of $0.005$ over $[2.30,2.50]$, the spread-minimising base
is $\beta = 2.400$ exactly. There is no detectable polynomial correction: the law is
$A(r) = c\beta^r$ with a single constant.

The tightest surviving pure-exponential form is
\[
   \text{(T1)}\qquad K \text{ D-reducible} \;\Longrightarrow\; a \ge 2.4^{\,r}/12.79,
\]
with $a \ge 2.4^r/12.78$ \textsc{killed}. Its per-ring margins lie in
$[1.001, 1.092]$; it is the only candidate near-tight at both ends and in the middle
at once. \textbf{(T1) is corpus-supported only.}

\begin{table}[t]
\centering\small
\caption{$f(r)$ implied by pairing a threshold with a cap. ``$\gamma=4/3$'' is the
sharp cap of Conjecture~\ref{conj:mass}, not proven.}
\label{tab:implied}
\begin{tabular}{lrrrrrrrrrr}
\toprule
$r$ & 8 & 9 & 10 & 11 & 12 & 13 & 14 & 15 & 16 & 20 \\
\midrule
\textbf{true $f(r)$} & \textbf{5} & \textbf{5} & \textbf{6} & \textbf{7} & \textbf{7} & \textbf{8} & ? & ? & ? & ? \\
weak cap $2^{r+k-3}$ (proven) & 2 & 2 & 2 & 3 & 3 & 3 & 4 & 4 & 4 & 5 \\
Thm~\ref{thm:star} $11\cdot2^{r+k-7}$ (proven) & 2 & 3 & 3 & 3 & 4 & 4 & 4 & 4 & 5 & 6 \\
$\gamma=3/2$ (proven at $k=2$ only) & 3 & 4 & 4 & 5 & 5 & 5 & 6 & 6 & 7 & 9 \\
$\gamma=4/3$ (the target) & 4 & 5 & 5 & 6 & 6 & 7 & 8 & 8 & 9 & 12 \\
\bottomrule
\end{tabular}
\end{table}

Table~\ref{tab:implied} shows the payoff. Pairing (T1) with the sharp cap gives
$f(r) \ge \lceil 0.63375\,r - 1.6306\rceil$, i.e.\ $4,5,5,6,6,7$ against the true
$5,5,6,7,7,8$: \textbf{off by at most one at every ring, from two inequalities}. That
is an analytic near-reproduction of the entire $f(r)$ theorem, valid for all $r$.
With the \emph{proven} cap the implied slope is only $0.263$, so the entire remaining
factor of $2.4$ in $f$-slope lives in the cap, not the threshold.

\begin{remark}[The structural principle]
Under a cap $B(r,k) = C\cdot 2^r\gamma^k$, the implied $f$-slope is
$(\log_2\beta - 1)/\log_2\gamma$ --- $0.263$ at $\gamma=2$, $0.450$ at
$\gamma=3/2$, $0.634$ at $\gamma=4/3$. \emph{The entire source of $f$-growth is the
excess of the threshold base $\beta$ over $2$}; a law with $\beta \le 2$ proves nothing
about $f$. This kills the whole wheel-anchored family of threshold candidates as a
route to $f$, since their $r$-dependence $2^r$ cancels against the cap's.
\end{remark}

\begin{remark}[Honest caveat on the pairing]
(T1)'s implied slope $0.634$ slightly \emph{exceeds} the empirical $f$-slope
($\approx 0.60$, the endpoint slope $(8{,}5)\to(13{,}8)$; least squares over the six
points gives $0.629$, which is not significantly below $0.634$). Extrapolated at the
endpoint slope the implied bound would overtake the truth around $r\approx 54$; at
the least-squares slope, not until $r\approx 263$. So at least one of \{(T1), the
sharp cap, the $0.60$ slope
estimate\} must weaken before then.
\end{remark}

\begin{remark}[Where the corpus does and does not constrain]
Only rings $6$--$10$ have machine-generated D-reducibles; at rings $11$--$16$ the
entire D-reducible population is the union of the three published catalogs, and at
rings $15$--$16$ there is not even a non-reducible negative population. A frontier
read off rings $14$--$16$ measures what catalog authors chose to publish. We tested
this: re-running the winners restricted to non-catalog sources ($51{,}067$ records
evaluated, of which $5{,}478$ are the D-reducible witnesses that can bind;
effectively rings $\le10$) leaves the generated-only frontier
$\{6{:}19, 7{:}39, 8{:}94, 9{:}228, 10{:}516\}$ versus the all-source
$\{6{:}16, 7{:}39, 8{:}94, 9{:}211, 10{:}496\}$, tightening the $\beta=2.4$ constant
by only $3.9\%$; and $a \ge 2.4^r/12$ is killed by the generated population
independently of any catalog. So the frontier is not a catalog artifact on the rings
where we can check --- but that check covers only $r\le10$.
\end{remark}

%======================================================================
\section{Methods: the verification architecture}
\label{sec:methods}
%======================================================================

The programme's operating assumption is that on this material a single
implementation's output is worthless, and that the interesting failures are the ones
a plausible-looking pipeline produces silently. Five mechanisms carry the load.

\paragraph{(M1) Two (and where needed three) independent implementations.} Our
reducibility checker is written from the published mathematics, not ported from the
released C. Where the two disagree, a \emph{third} implementation with no shared
code path adjudicates (\art{src/fourcolor/brute\_force.py}: different edge indexing,
different triangle enumeration, no gauge fixing, no greedy variable ordering). The
same discipline is applied to the discharging side: the Python port of
\texttt{always\_apply}/\texttt{never\_apply}/blocking was differentially validated
against the C++ on all $5{,}692{,}937$ wheel candidates before being trusted, and
re-validated on all $10{,}094$ published leaves before being used in
\S\ref{sec:lp}. In both places the differential test found real defects --- in the
legacy C in one direction (\S\ref{sec:verify}) and in our own serializer in the other
(\S\ref{sec:fr-enum}).

\paragraph{(M2) Adversarial mutation testing.} A candidate statement that holds on a
pool has, on the evidence of \S\ref{sec:negative}, essentially no chance of being
true. \art{tools/stress\_test\_theorem.py} attacks a candidate by generating mutants
of pooled configurations under a single diagonal edge flip
(\art{src/fourcolor/mutate.py}), filtering each mutant through a structural legality
gate (a port of RSST's own \texttt{ReadConf} conditions), then through the candidate
rule on recomputed features, then through a canonical-key novelty check; survivors
are labelled by the reducibility checker and every false positive is re-adjudicated
by the compiled C oracle. The one time we ran this on genuinely novel candidates it
killed all three of them (\S\ref{sec:negative}).

\paragraph{(M3) The lemma harness: test before proving.} Any candidate statement can
be machine-checked against the entire labelled corpus of $59{,}142$ canonically
deduplicated configurations (rings $6$--$16$, $13{,}169$ of them D-reducible)
in about $0.2$ s for a scalar statement:
\begin{center}\small
\texttt{tools/test\_lemma.py -{}-check "rec.d\_reducible implies rec.a >= 2.4**rec.r/12.79"}
\end{center}
Each run appends a receipt to \art{results/mass-law/lemma\_log.jsonl}, content-hashed
against the corpus signature \texttt{a2b1feba331bf2962b2b} --- a blake2b digest over
\texttt{(canonical\_hash, r, n, a, b, d\_reducible)} for records in canonical-hash
order, so a receipt identifies exactly which configurations with which labels a
verdict was taken over. Every corpus-supported claim in \S\ref{sec:mass} carries such
a receipt. The harness's verdict vocabulary is deliberately restricted: \textsc{holds}
means ``survived the corpus'', never ``proved''. Of the $87$ logged receipts, $66$ are
\textsc{holds} and $21$ are \textsc{killed}; the killed ones are as much of the
output as the surviving ones.

\paragraph{(M4) Exact-arithmetic certificates with re-derivation.} Where a result is
a computation, we ask it to emit an object that an independent, simpler program can
check. \art{tools/verify\_farkas.py} is the archetype: solver-free,
\texttt{fractions.Fraction} throughout, and --- the part that matters --- it does not
trust the certificate's stored coefficients, re-deriving every row from its source
artifact and failing loudly on mismatch. Corrupting one coefficient by hand is
rejected with an exact claimed-versus-re-derived diff. The corresponding discipline
on the enumeration side is that generation completeness is delegated to a pinned
canonical-construction-path generator \cite{brinkmann2007} rather than to our own
code, and legality to a port of the original authors' own parser.

\paragraph{(M5) Adversarial review of every claim before it is written down.} Each
major result was handed to a reviewer prompted to \emph{refute} it and to default to
``refuted'' when unsure, with independent reimplementation required. This produced,
among other things: the discovery that the pool-independence claim in
Lemma~\ref{lem:mono} was false (fixed); the code-necessary versus math-necessary
distinction that became \S\ref{sec:necessity}; the coarse-cartwheel gap in the
low-degree rows (closed exactly, \S\ref{sec:lowdeg}); the demotion of a claimed
$\gamma=3/2$ cap for all $k$ to Theorem~\ref{thm:gamma32} at $k=2$ only, after the
reviewer showed the enabling hypothesis fails on $90.5\%$ of the corpus; and two
false provenance sentences about \cite{nl4ct2026} that had survived several drafts.
We record this as a method because in every one of those cases the error was in
prose or scope rather than in a computation, and no amount of differential testing
would have caught it.

\paragraph{Test suite and pinning.} $23$ test files, $295$ discovered test cases,
run by \texttt{make test}; at the time of writing $293$ pass and $2$ fail. Both
failures are a data-drift guard firing correctly rather than a defect: the ring-$8$
trace corpus grew from $1{,}162$ to $1{,}736$ records after the ring-$8$ model was
trained, so \texttt{verify\_split\_against\_metrics} refuses to reconstruct that
model's train/validation split against the enlarged data. Ring $8$ carries none of
the results reported in \S\ref{sec:ml}, which are at rings $9$ and $10$. All third-party inputs are pinned: four arXiv e-prints
with recorded SHA-256 digests (\art{third\_party/CHECKSUMS.sha256}), three git
commits for the released 2026 artifacts (\art{third\_party/GIT\_PINS.txt}), and
\texttt{plantri} 5.5 with the digests of both binary and source.

\section{Machine learning: what was tried, and what it gave}
\label{sec:ml}

This section is an accounting, not an advertisement. None of the theorems in
this paper were discovered by a neural model. They came from exhaustive
enumeration, LP duality, and machine-checked mining over finite catalogs, as
described in the preceding sections. We nonetheless ran a substantial
interpretability programme alongside that work, on the hypothesis that a
small model trained to predict D-reducibility might expose a human-statable
sufficient condition that the enumeration was not organized to find. We
report what that programme gave: one modest positive result, one
methodological finding about the fragility of ``mechanism'' as a concept,
and several negative results we think are worth recording so that others do
not have to rediscover them. We write in the register of a referee report
on our own work.

\subsection{Design and motivation}

The programme deliberately transplanted the template of
Huang, Jackson, and Lee's mechanistic-interpretability study of the zeta
map \cite{huang2025zeta}: train a small transformer on a well-defined
mathematical computation, probe its internal representations, and try to
extract a human-readable mechanism. That template sits inside the broader
agenda of using machine learning to guide mathematical intuition
\cite{davies2021}. The target computation here was D-reducibility itself:
given a configuration, predict the Kempe-closure trace produced by the
D-reduction search and the final verdict (reducible or not). The rationale
was simple. If a small model can predict reducibility reliably, whatever
internal feature it relies on is a candidate for a sufficient condition a
human could then state and verify independently -- the same logic that
motivated the zeta-map study's search for a legible mechanism.

We built two model generations.

\textbf{D1.} An encoder-decoder transformer ($d_{\mathrm{model}}=192$,
4 heads, feed-forward width 384, 2 encoder layers, 2 decoder layers,
dropout 0, 1{,}543{,}725 parameters), trained seq2seq from a canonical
token sequence encoding a configuration to a target sequence of
closure-trace count buckets followed by a verdict token. Training ran 120
epochs at learning rate $10^{-3}$, batch size 64, seed 0, for 10{,}392~s on
an Apple MPS device. The corpus, \art{data/d1\_corpus.jsonl}, held
14{,}036 deduplicated records spanning ring sizes 6--16 (633 from RSST,
2{,}614 from Steinberger, 5{,}961 generated, 4{,}828 drawn from the 2026
pool); after excluding records with degenerate traces or out-of-range
length, 6{,}559 remained, split 5{,}904 train / 655 validation. Artifacts:
\art{checkpoints/d1\_full\_v1.pt},
\art{results/d1/train\_metrics\_full\_v1.json},
\art{results/d1/corpus\_stats.json}.

\textbf{D1-v2.} A second generation of per-ring specialists at
$r = 8, 9, 10$, sharing a 2-layer encoder ($d_{\mathrm{model}}=256$, 4
heads, feed-forward width 512, masked-mean-pool summary) with two
alternative heads. A ``dynamics'' head iterates a weight-tied
\texttt{GRUCell} over the per-ring survivor-set bitmap with per-round
supervision (3.26M parameters at $r=9$, 6.64M at $r=10$); a
``verdict-only'' \textsc{control} head reads the verdict directly off the
encoder summary with no recurrence (1.19M parameters), and exists
specifically to test whether the dynamics head's recurrent machinery
contributes anything the summary vector does not already contain. Data
lived in \art{data/v2/traces\_r8.jsonl}, \art{data/v2/traces\_r9.jsonl},
and \art{data/v2/traces\_r10.jsonl}, carrying full per-round survivor sets;
the corpora held 1{,}162 / 1{,}446 / 1{,}746 configurations at $r=8,9,10$, with
ring-colouring code-index sizes 1{,}094 / 3{,}281 / 9{,}842 respectively. Artifacts: \art{results/d1v2/r9/},
\art{results/d1v2/r10/}, \art{results/d1v2/r10v2/}.

\subsection{Results}

D1 reaches verdict accuracy 0.9298 on its 655-example validation split,
against a majority-class baseline of 0.5603 -- a real signal, well above
chance. Full-sequence exact match (verdict token plus the entire
closure-trace bucket sequence) is a much weaker 0.2397: the model gets the
answer right far more often than it gets the reasoning trace right.

We then asked what the encoder actually represents, using a probe ladder
recorded in \art{results/d1/interp/report.md} (ridge logistic probes on
standardized inputs, probe-test $n=131$, 95\% bootstrap confidence
intervals with \texttt{n\_boot}$=2000$). A 16-feature hand-coded shallow
baseline -- counts and simple graph statistics of the configuration, no
learned representation involved -- already reaches 0.8779
$[0.8168, 0.9313]$. The best pure-encoder representation is
\texttt{encoder\_layer0\_mean}, the \emph{pre-attention} embedding bag, at
0.9008 $[0.8473, 0.9466]$: a delta of $+0.0229$ over the shallow baseline,
but with overlapping confidence intervals, so not a distinguishable
improvement. The post-attention layers are worse, not better: layers 1 and
2 both give 0.8855, despite fitting the training set considerably better
than layer 0 (train accuracy 0.9599 and 0.9866 respectively, against 0.8817
for the pre-attention bag). Attention is improving the model's memorization
of the training distribution without improving the generalizing signal a
downstream probe can extract. Augmenting the shallow features with
hand-derived structural candidates reaches 0.9237
$[0.8779, 0.9695]$ -- that is, hand features beat the encoder's best
representation outright. We state this plainly as the honest summary of
D1: the transformer's internal representation of reducibility is, on this
evidence, no better than what a referee could write down with a pencil,
and in places worse.

A second question was whether the decoder accumulates verdict evidence as
it emits the closure trace, the way one might hope a legible mechanism
would. We probed decoder hidden states at each greedy decoding step
$t = 0, \dots, 39$. Verdict-probe accuracy is 0.8473 at $t=0$, before a
single trace token has been produced, and the curve stays essentially flat
thereafter: the mean over $t=0$--$2$ is 0.8702, the mean over $t=35$--$39$
is 0.8656, and every later confidence interval overlaps the one at $t=0$.
The closure-trace tokens ride along with a verdict decision already
implicit at the start; they do not build toward one.

\subsection{The boundary result}
\label{sec:ml-boundary}

The results above concern accuracy on the full validation distribution,
which is dominated by easy cases. We pre-registered a harder ``boundary''
stratum, flagged by \texttt{tools/d1v2\_datagen.py::boundary\_flag}, to test
whether any representation carries information beyond the shallow features
specifically on the cases an easy classifier would get wrong: a
configuration is boundary if it is non-reducible with
$0 < n_{\mathrm{consistent}} \le 50$, or reducible with rounds $\ge 6$. We
pre-registered the decision rule before looking at boundary results: a
representation beats the shallow baseline only if its boundary 95\%
bootstrap confidence interval does not overlap the shallow baseline's and
lies strictly above it.

The first pass, at rings 9 and 10 on the smaller corpus, was a clean null:
0 of the 7 non-baseline representations at $r=9$ and 0 of 7 at $r=10$ cleared the
rule, and
the boundary strata in the probe-test splits were tiny (7 and 11 examples
respectively) -- underpowered by construction.

The second pass used an enlarged ring-10 corpus (3{,}269 configurations;
probe-test $n=98$, of which 21 are boundary) and gave a real result.
Table~\ref{tab:boundary} summarizes it.

\begin{table}[htbp]
\centering
\begin{tabular}{lcc}
\toprule
Representation & Overall accuracy & Boundary accuracy [95\% CI] \\
\midrule
shallow & 0.7551 & 0.5238 $[0.3321, 0.7143]$ \\
shallow\_plus\_structural & 0.8571 & 0.7619 $[0.5714, 0.9048]$ \\
control\_encoder\_summary & 0.9082 & 0.9048 $[0.7619, 1.0000]$ \\
dynamics\_encoder\_summary & -- & 0.9048 \\
dynamics\_h0 & -- & 0.9048 \\
dynamics\_h2 & -- & 0.9048 \\
\bottomrule
\end{tabular}
\caption{Ring-10 boundary probe results, second pass ($n=98$ overall,
21 boundary). \texttt{control\_encoder\_summary} and three dynamics-model
representations clear the pre-registered rule against the shallow
baseline.}
\label{tab:boundary}
\end{table}

\texttt{control\_encoder\_summary} -- the 256-dimensional encoder summary
feeding the verdict-only control head, with no recurrent dynamics involved
at all -- clears the pre-registered rule against the plain shallow
baseline, as do three dynamics-model representations. This is the one
genuinely positive result in the programme: some encoder representation
carries boundary-relevant information the 16 hand-picked shallow features
do not.

We caveat it on three counts. First, $n=21$ boundary examples is a small
stratum; the confidence intervals in Table~\ref{tab:boundary} are wide.
Second, against the also-hand-engineered \texttt{shallow\_plus\_structural}
baseline the comparison is never confidence-interval-decisive at any seed
we tried (0.9048 versus 0.7619, intervals overlap); the result holds
against the naive 16-feature baseline, not against a reasonably strengthened
one. Third, we note for the record that the source report's own concluding
prose (\art{results/d1v2/v2run/interrogation.md}) contradicts its table: it
recycles wording from the earlier null run at rings 9 and 10. We report the
table, not the prose.

\subsection{Extraction: hand features recover most, but not all, of the advantage}
\label{sec:ml-extraction}

Given that some encoder representation beats the shallow baseline on
boundary cases, we tried to extract what it was finding, in the spirit of
turning a learned feature into a stated one. We hand-coded 14 arrangement
features of $H_5$, the subgraph induced on interior vertices of degree
exactly 5 (\art{src/fourcolor/d1\_features.py}): number of connected
components, largest component size, number of isolated vertices, edge
count, density, maximum degree, existence of a length-3 path, number of
vertices adjacent to the ring, minimum and mean component-to-ring
distance, minimum and mean inter-component distance, the count of
degree-5 vertices with at least 3 degree-5 neighbours, and the existence
of a diamond in $H_5$.
Appending these 14 to the 24-dimensional \texttt{shallow\_plus\_structural} vector
(16 shallow features plus 8 structural ones) gives a 38-dimensional
\texttt{shallow\_plus\_arrangement} vector. Table~\ref{tab:extraction}
reports boundary accuracy across four seeds, drawn from
\art{results/d1v2/extraction/report.md}.

\begin{table}[htbp]
\centering
\begin{tabular}{lccccc}
\toprule
Seed & $n_{\mathrm{boundary}}$ & shallow & +structural & +arrangement [95\% CI] & control\_encoder \\
\midrule
0 & 21 & 0.5238 & 0.7619 & 0.9524 $[0.8571, 1.0000]$ & 0.9048 \\
1 & 23 & 0.4348 & 0.7826 & 0.9130 $[0.7826, 1.0000]$ & 0.7826 \\
2 & 24 & 0.5417 & 0.7083 & 0.8750 $[0.7083, 1.0000]$ & 0.7500 \\
3 & 25 & 0.5200 & 0.7600 & 0.7600 $[0.5990, 0.9200]$ & 0.9200 \\
\bottomrule
\end{tabular}
\caption{Boundary accuracy of the hand-coded arrangement-feature extraction
against the plain shallow baseline, the structural baseline, and the
encoder control, across four seeds.}
\label{tab:extraction}
\end{table}

The extraction succeeds on three of four folds: at seeds 0, 1, and 2 the
arrangement features push boundary accuracy well above
\texttt{shallow\_plus\_structural}, and in all three above the encoder control as
well (0.9524 vs 0.9048, 0.9130 vs 0.7826, 0.8750 vs 0.7500). It fails on seed 3: the arrangement features add
exactly nothing over the structural baseline there, an identical point
estimate of 0.7600 to 0.7600, while the encoder control reaches 0.9200. We
do not know why seed 3 is different; we simply report that it is.

The boundary result of Table~\ref{tab:boundary} itself replicates across
these same four seeds in direction but not always in confidence-interval
strength: the control encoder beats the plain shallow baseline by $+0.38$,
$+0.35$, $+0.21$, and $+0.40$ at seeds 0--3 respectively, so the sign is
consistent 4 of 4. But the confidence-interval-decisive version of the
comparison -- the pre-registered rule -- holds only at seeds 0 and 3, i.e.\
2 of 4. The headline boundary result reported in
of \S\ref{sec:ml-boundary} was therefore somewhat fortunate in its seed choice;
a differently seeded first pass could plausibly have reported an
inconclusive result instead. Artifacts:
\art{results/d1v2/extraction/report.md},
\art{results/d1v2/replication/report.md}.

We also performed a small error-set digestion at seed 0
(\art{results/d1v2/digestion/cases.md}). Of the 21 boundary probe-test
examples, the encoder control is right and the structural baseline is
wrong on 4; the reverse holds on 1; both are right on 15; both are wrong
on 1. All 4 configurations where the encoder wins have $n=18$, are
D-reducible, and have closure round counts $6, 6, 7, 6$. Their mean count
of degree-5 diamonds is 2.75, against 0.88 for the rest of the validation
boundary set and 0.97 corpus-wide -- roughly a factor of three -- with
every other feature we measured flat across the two groups. We decline to
run a significance test at $n=4$; the sample is too small to support one
honestly.

We flag one honest gap here. The automated per-configuration check that
the arrangement features specifically fix these four cases did not
actually run: \art{tools/d1v2\_extraction\_test.py} guards its
\texttt{error\_set\_check} behind
\texttt{hasattr(interro, "probe\_predictions")}, which is \texttt{False}
for the interrogation object actually produced, so the check is silently
skipped and \texttt{error\_set\_check} is empty at every seed. The causal
link we would like to draw -- diamond count explains the four flipped
cases -- is therefore suggestive, resting on the aggregate statistic above,
and not established at the level of individual predictions.

\subsection{Mechanism is training-contingent: a replication finding}

This subsection reports a companion experiment run in a separate
repository, \texttt{/Users/yugendren/experiments/zeta\_map\_interp}, not in
the codebase for this paper. Its subject is the published zeta-map result
we took as our template, \cite{huang2025zeta}, and the question is whether
its reported mechanism replicates under a faithful reimplementation.

We reimplemented the published architecture (1 encoder layer, 1 decoder
layer, 1 attention head, $d_{\mathrm{model}}=128$, feed-forward width 256,
post-norm, learned positional embeddings, weight-tied output projection):
337{,}029 parameters against the paper's reported 339{,}716. Training
reached the same headline accuracy the paper reports: 0.9950 greedy
exact-match on a held-out instances of size $n=13$. We then reran each of the
paper's three reported mechanistic findings against our own checkpoint.

\begin{itemize}
\item \emph{Causal ablation.} The paper reports that ablating
cross-attention to up-steps costs about 1.5\% exact-match accuracy, read
as evidence of a selective mechanism attending preferentially to up-steps.
We measure a baseline of 0.996 exact match falling to 0.0000 when
ablating attention to up-steps, and to 0.0005 when ablating attention to
the complementary down-steps. Both ablations are catastrophic and
symmetric: cross-attention is load-bearing wholesale in our model, with no
selectivity toward either step type.

\item \emph{Attention statistics.} Mean cross-attention mass on up-steps
is 0.5344 against a uniform baseline of 0.5; the attention argmax lands on
an up-step for 0.5739 of queries; only 0.0682 of queries have their argmax
at the path's own maximum level; the mean level of the argmax-attended
position is 2.369, against 2.471 averaged over all source positions. These
are essentially uniform statistics -- we find no evidence of attention
selecting the path's highest level.

\item \emph{Linear probes.} The paper reports approximately 84\% probe
accuracy for recovering path levels from encoder hidden states. We measure
0.9478 (14 classes, majority baseline 0.2439) for level, and 1.0000 for
step type (2 classes, baseline 0.5). This one direction replicates, and in
fact strengthens.

\item \emph{Length generalization.} The paper evaluates lengths
$n = 11$--$16$. Our model scores exactly 0.0000 exact match at
$n = 11, 12, 14, 15, 16$, with per-token character accuracy 0.5872 at
$n=12$ against roughly 0.5 chance. We tried four positional-placement
variants at evaluation time (left-aligned to length 24, left-aligned
padded to length 26, right-aligned, centered); all four give 0.0000 exact
match, which rules out a pure positional-embedding-range explanation for
the failure.
\end{itemize}

We draw the conclusion carefully, and without asserting that the original
authors made an error: same architecture, same task, the same held-out
accuracy we would need to call the reimplementation faithful, and yet a
materially different internal mechanism by every measure we could apply.
Whatever mechanism a small transformer settles on for a mathematical
computation appears to be, at least in part, a property of the particular
training run rather than a property of the task it was trained on. For an
interpretability-driven discovery programme this is the central
methodological risk: an extracted ``mechanism'' may describe a training
run rather than a fact about the mathematics. This is exactly the failure
mode we then observed inside our own setting, at seed 3 of
\S\ref{sec:ml-extraction}.

\subsection{Assessment}

None of the four contributions of this paper came from the neural models.
The one thing the machine-learning line delivered that fed forward into
the rest of the paper was the degree-5 arrangement feature family of
\S\ref{sec:ml-extraction}, which sharpened our attention on $H_5$ and became
one input, among others, to the rule-mining process that produced the
candidate rules discussed in Section~\ref{sec:negative} -- of which the
surviving rule, in its exhaustively verified form, became
Theorem~\ref{thm:fr}. We describe that as an indirect, heuristic
contribution and make no stronger claim than that.

We also record two structural failures of the programme, so that the
positive results above are read in proportion. First, the dynamics
model -- the entire interpretability design bet, and the closest analogue
to the zeta-map paper's step-by-step mechanism -- attains teacher-forced
verdict accuracy 1.0000, but collapses under free-running rollout to
0.7778 at $r=9$, 0.7786 at $r=10$, and 0.6245 at $r=10$ in the second
($r10v2$) configuration, with free-running boundary accuracy exactly
0.0000 at $r=10$ and $r10v2$: unrolled without teacher forcing, the model
converges to the trivial fixed point of predicting that nothing is ever
eliminated. Second, the D1 checkpoint's learned positional embeddings
extend only to length 259, which admits 7{,}021 of the 14{,}036
configurations in the corpus; ring sizes above roughly 13 were structurally
out of the model's reach, independent of anything about the underlying
mathematics.

Finally, we caught and quarantined one tautological feature during this
work. The quantity $n_{\mathrm{consistent}}$ correlates with the verdict
label at Spearman $-0.850$ and appears to ``close 167\% of the gap'' to
perfect boundary accuracy -- because, on inspection, it is the label in
disguise: $n_{\mathrm{consistent}} \le 50$ is one of the two disjuncts
defining boundary non-reducibility in the first place. It is excluded from
consideration via \texttt{TAUTOLOGICAL\_CANDIDATES} in
\art{tools/d1\_interrogate.py}, and we mention it here only as a caution to anyone repeating this
kind of search: a feature that predicts the label unusually well is worth
checking for circularity before it is reported as a discovery.


%======================================================================
\section{Negative results and killed claims}
\label{sec:negative}
%======================================================================

We report these at length because the verification discipline of
\S\ref{sec:methods} is only credible if its failures are visible, and because
each of the three killed claims below was, at the moment it was killed, the most
attractive result the project had.

\subsection{Three mined sufficient conditions, killed by mutation}

We mined a pool of $16{,}794$ canonically deduplicated configurations
($11{,}919$ D-reducible / $4{,}875$ not, rings $6$--$16$, $0$ label conflicts) for
minimal conjunctions of structural predicates that are sufficient for
D-reducibility with zero false positives on the pool. The search space contains
$81{,}386$ such minimal zero-false-positive rules. The three best by coverage were:

\begin{center}\small
\begin{tabular}{clrr}
\toprule
\# & rule & coverage & pool FPs \\
\midrule
1 & $n \ge 25$ \textsc{and} $\overline{\deg}_{\rm int} \ge 5.7\overline{3}$ \textsc{and} $h_5\text{-density} \le 2/7$ & 2{,}313 & \textbf{0} \\
2 & $n \ge 25$ \textsc{and} max run of constant ring degree $\le 4$ \textsc{and} $n_{\deg 5\to\ge3\deg 5} = 0$ & 2{,}240 & \textbf{0} \\
3 & $k \ge 11$ \textsc{and} $h_5\text{-density} \le 2/7$ \textsc{and} $n_{\deg 5\to\ge3\deg 5} = 0$ & 2{,}234 & \textbf{0} \\
\bottomrule
\end{tabular}
\end{center}

All three survived $5$-fold cross-validation, but that survival was already weak
evidence: across the five folds their coverage ranks are $1,1,\text{subsumed},38,1$
for rule 1, $2,2,\text{subsumed},47,7$ for rule 2 and $12,4,\text{subsumed},43,4$
for rule 3 --- rule 3 never reaches first or second in any fold, and all three fall
far down in fold 4. All three are \textbf{false}.

Adversarial mutation search (\S\ref{sec:methods}, M2) generated $339{,}600$ mutants
from $9{,}391$ seeds; $248{,}861$ failed the structural legality gate, $88{,}496$ no
longer satisfied the rule, $291$ were duplicates, and $1{,}952$ novel legal
rule-satisfying configurations were submitted to the reducibility checker. Of those,
\textbf{$236$ are false positives} --- configurations satisfying the rule that are
not D-reducible: $83$ for rule 1, $44$ for rule 2, $109$ for rule 3, all with
distinct canonical keys, at rings $13$ ($54$), $14$ ($105$), $15$ ($56$) and $16$
($21$), arising from $113$ distinct seed configurations. The false-positive rate
among labelled mutants is $12.1\%$. \textbf{$215$ of the $236$ were independently
confirmed} by the compiled C oracle, with $0$ disagreements between the two
implementations; the remaining $21$ --- exactly all the ring-$16$ cases --- could not
be adjudicated because the oracle's output truncates before its verdict line, the
same failure mode as \S\ref{sec:verify}.

\begin{quote}
\textbf{The lesson, stated as we now apply it:} zero false positives on a pool of
$16{,}794$ configurations, stable across cross-validation folds, is not evidence of
soundness. It is evidence that the pool does not contain the counterexamples, which
in a curated catalog is what one should expect. Always mutate-stress.
\end{quote}

\subsection{The rule that survived --- and why the mutation test did not really test it}

The dual, \emph{necessary}-direction rule --- ``$r=11$ and $k \le 6$ implies
not D-reducible'', covering $1{,}630$ non-reducible configurations with zero false
positives --- was the one candidate the mutation search did not kill, and it is the
statement that became the $r=11$ case of Theorem~\ref{thm:fr}.

Its ``survival'' should nevertheless be discounted almost entirely. Of $41{,}449$
mutants generated, $41{,}148$ failed the legality gate, $273$ failed the rule and
$28$ were duplicates: \textbf{zero novel rule-satisfying mutants were ever submitted
to the checker}, and the run took one second. The tool's verdict is computed
mechanically as ``killed if any false positive, else survived'', so
\textsc{survived} here means ``no counterexample was found'', not ``counterexamples
were sought and none exist''. The reason is structural and instructive: configurations
at $r=11$ with $k \le 6$ sit on the legality boundary --- their interior vertices have
degree exactly $5$ --- so almost every single edge flip drops a vertex below its
minimum degree. Raising the per-seed budget to $400$ tries still yielded nothing
testable.

The real evidence for that statement is therefore not the mutation test but the
exhaustive \texttt{plantri} sweep of \S\ref{sec:fr-enum}. We flag this because
presenting the mutation result as independent confirmation would be exactly the kind
of thing this section exists to prevent.

\subsection{The strong LP encoding, killed by its own negative control}

The natural formalisation of ``the discharging closes'' --- \emph{every unblocked
stage-1 wheel must end with strictly negative charge} --- is wrong, because the 2026
schema permits survivors and the published proof leaves $16{,}148$ stage-1 wheels
alive on purpose. We implemented it anyway, with the max-of-affine cutting-plane
machinery it requires, and ran the \emph{full-pool} control first. The control is
\textsc{infeasible} at iteration $0$ with a verified $10$-row certificate
($y^\top b = -175$): the encoding ``disproves'' the published theorem. We discarded
the encoding on that basis and built the leaf-level one of \S\ref{sec:lp}
instead.

We stress the ordering. The control was run \emph{before} the ring-$\le 14$ variant,
and the ring-$\le14$ variant under that encoding is also infeasible --- with a nicely
large $y^\top b = -380$. Had we run the two in the other order and stopped, we would
have published a false result with a valid-looking certificate. A negative control
that can fail is the only defence against this, and it is cheap.

\subsection{The naive threshold law, killed by the harness}

The first form of the threshold half of Conjecture~\ref{conj:mass}, read off the
$r=8$ minimum and the observed growth ratio, was
$K \text{ D-reducible} \Rightarrow a \ge 94 \cdot 2.4^{\,r-8}$. The harness kills it
with $43$ counterexamples in about $0.2$ s. Also killed: $a \ge 2.4^{\,r}/12$
($15$ counterexamples, and still $2$ counterexamples when restricted to the
machine-generated population, so this is not a catalog artifact);
$a \ge 5\cdot2^{\,r-4}$ ($6$); $a \ge 2^{\,r-1}$ ($311$); $a \ge k\cdot 2^{\,r-4}$
($398$); $a \ge W(r)\cdot 1.3^{\,k-1}$ with no leading constant (all $13{,}169$);
$a \ge 3^{\,r-1}/2/(0.25\,r^2)$ ($17$). On the cap side the harness killed
$a \le 5\cdot2^{\,r+k-6}$ with exactly six counterexamples --- the six wheels --- which
is precisely what Theorem~\ref{thm:star}'s $\min$-over-$h$ form predicts, and is the
reason the unconditional constant is $11$ and not $10$. The surviving forms in
\S\ref{sec:mass} are what is left after that attrition.

\subsection{Two further honest notes}

\paragraph{A conjecture we would have overclaimed.} An earlier draft of
\S\ref{sec:mass} stated $\gamma=3/2$ for all $k$, under a hypothesis (H2) about arc
lengths that looked innocuous. (H2) fails on at least $53{,}549$ of $59{,}142$
records ($\ge 90.5\%$) and on $37.66\%$ of individual BFS steps ($140{,}084$ of
$372{,}008$). The adversarial reviewer called
presenting the conclusion under that hypothesis ``a real overclaim'', and it was
demoted to Theorem~\ref{thm:gamma32}.

\paragraph{Corrections to our own record.} Two numbers circulated in this project's
internal notes that are wrong and that we correct here: the mutation search produced
$236$ false positives, not $26$ (an apparent dropped digit, propagated through three
documents); and the wheel-candidate total at hub degrees $7$--$11$ is $5{,}692{,}937$,
not $5{,}690{,}937$. Both are corrected in this paper and neither affects a verdict.

%======================================================================
\section{Discussion and open problems}
\label{sec:open}
%======================================================================

\paragraph{The sharp cap via transfer matrices.} Theorem~\ref{thm:obstruction} says
the sharp cap cannot be reached through $P(S,4)$; one needs a direct handle on
$|\Phi(K)|$. Two routes remain. (a) An injection
$\Phi(K) \hookrightarrow \Phi(K')\times[4/3]$ for a configuration $K'$ with one fewer
interior vertex \emph{and the same ring}. Ordinary vertex deletion will not do:
deleting an interior vertex of degree $d$ with $j$ consecutive ring neighbours yields
ring length $r+d-2j+2$, so the ring \emph{grows} --- for $d=5$, $j=2$ it grows by $3$,
multiplying the cap by $8$, and the induction would have to supply
$a(K)\le a(K')/6$, an inequality of the wrong strength. Deletion-based induction on
this class is a dead end unless the ring can be kept fixed. (b) A transfer-matrix
bound on the \emph{ring-trace} operator, whose state space is ring colourings rather
than tri-colourings. The Jacobsen--Salas--Sokal Temperley--Lieb machinery
\cite{jacobsen2006} is the right formalism, and its state space --- non-crossing
partitions of the ring --- is exactly the object $\Phi$ lives on. We regard this as
the highest-value open problem in the programme. Failing that, proving $\gamma=3/2$
in general (H2 is a finite, checkable link condition) already lifts the implied
$f$-slope from $0.263$ to $0.450$ and is permitted by the Bridge route.

\paragraph{Removing the fixed-shapes caveat.} Theorem~\ref{thm:lp} fixes the $84$
published rule shapes. The natural strengthening --- ``no discharging schema of
bounded rule complexity closes over the ring-$\le14$ pool'' --- requires column
generation: given a violated leaf, price out a new planar degree-interval pattern
that would fix it. The pricing subproblem is a search over planar patterns with a
charge dart, not an LP, and it is the main technical obstacle. A tractable
intermediate is to enumerate all rule shapes below a fixed size bound and re-run the
same certificate machinery over the enlarged column set: infeasibility there would be
a substantially stronger statement, and the certificate format is unchanged.

\paragraph{$f(r)$ beyond $r=13$.} The $r=13$ interior-$7$ cell alone required
$11{,}534{,}001$ plantri graphs and $\approx 30$ core-hours, $8$-way sharded
($\approx 3.7$ h wall); $r=14$ at
interior $8$ is roughly two orders of magnitude larger and is not reachable by the
present pipeline on a workstation. Three ways forward: (i) prune the generation with
a necessary condition on $a$ --- the mass law, if proven even in weak form, would
eliminate most cells without running the reducibility check; (ii) exploit the
observation of \S\ref{sec:mass} that the minimum-$a$ and minimum-interior frontiers
coincide, to search only near the frontier; (iii) accept a conditional result. Note
that the corpus already gives $f(14)\le 9$, $f(15)\le 10$, $f(16)\le 10$ from catalog
witnesses.

\paragraph{Formalisation.} The finite content of Theorem~\ref{thm:fr} is a replay of
$10{,}319$ configuration checks. That is within reach of a proof assistant given
Gonthier's formalisation \cite{gonthier2008} of the reducibility machinery, and it
would upgrade the theorem from ``certified'' to ``proven'' in the strongest available
sense. The Farkas certificate of Theorem~\ref{thm:lp} is even more tractable: $12$
rows over the rationals, though the row-derivation step (from cartwheel to
coefficients) carries the real formalisation burden.

\paragraph{The set-size question.} Our original motivation was to minimise the number
of configurations in an unavoidable set. We deliberately do not report a headline
number, because a minimal $U$ in the wheel-based schema is \emph{not} comparable to
RSST's $591$/$633$ (axle-based, C-reducibility with contracts, ring $\le14$). Any such
claim must be schema-qualified. Theorem~\ref{thm:lp} is the useful residue of that
line: an impossibility rather than an improvement.

\bibliographystyle{plain}
\bibliography{refs}

\appendix

%======================================================================
\section{Claim-to-artifact table}
\label{sec:artifacts}
%======================================================================

Every numerical claim in the paper is listed here with the repository artifact that
produced it. Paths are relative to the repository root. A claim marked
\textsc{corpus} is a lemma-harness verdict over the $59{,}142$-configuration
labelled corpus with signature \texttt{a2b1feba331bf2962b2b} and carries a receipt in
\art{results/mass-law/lemma\_log.jsonl}; \textsc{corpus} means ``survived'', never
``proved''.

{\small
\begin{longtable}{p{0.40\textwidth}p{0.52\textwidth}}
\toprule
\textbf{Claim} & \textbf{Artifact and content} \\
\midrule
\endfirsthead
\toprule
\textbf{Claim} & \textbf{Artifact and content} \\
\midrule
\endhead
\multicolumn{2}{l}{\emph{\S\ref{sec:lp}: the ring-$\le14$ obstruction}}\\[2pt]
$84$ rule shapes; $r(R)\in\{1,2\}$ &
  \art{third\_party/discharging-rules/R/} (84 \texttt{.rule} files);
  \art{third\_party/arxiv-2603.24880/src/main\_\_1\_\_\_1\_.tex} lines 624--625 \\
$8{,}200$-config pool; $5{,}895$ of ring $\le14$ &
  \art{third\_party/reducible-configurations/D/} (8{,}200 files);
  \art{build/steinberger-pool-r14/D} (5{,}895); ring histogram recomputed from
  the \texttt{N R} header of every file \\
Affine leaf form validated on $10{,}094$ published leaves, $0$ discrepancies &
  \art{tools/lp\_discharge.py -{}-validate}; leaves in
  \art{third\_party/computer-checks/wheels/zero/} \\
$8$-variant LP matrix (Table~\ref{tab:lpvariants}) &
  \art{results/steinberger/lp\_variants\_summary.json},
  \art{results/steinberger/lp\_control\_stats.json},
  \art{results/steinberger/lp\_main\_stats.json} \\
$19$-row Farkas certificate, $y^\top b=-30$ &
  \art{results/steinberger/lp\_main\_certificate.json} \\
$12$-row irreducible core, $y^\top b=-20$ &
  \art{results/steinberger/lp\_main\_certificate\_iis.json} \\
Exact-rational verification, rows re-derived &
  \art{tools/verify\_farkas.py} $\to$ \texttt{FARKAS CERTIFICATE VERIFIED}, exit 0 \\
All $12$ rows mathematically necessary &
  \art{results/steinberger/iis\_math\_necessity.json},
  \texttt{all\_mathematically\_necessary: true} \\
$11$ of $5{,}292$ (rule, out-dart) pairs undetermined &
  \art{tools/lp\_discharge.py -{}-audit-necessity} \\
Low-degree refinement counts (4 necklaces, up to $2.4\times10^8$) and the
  $3{,}046$-wheel sweep &
  \art{tools/lowdeg\_refinement\_check.py},
  \art{results/steinberger/lowdeg\_refinement\_check.json} \\
$580$ / $2{,}466$ unblocked $d{=}5$ / $d{=}6$ wheels, matching the C++ &
  released C++ \texttt{-{}-enum\_wheels -d 5}, \texttt{-d 6}; same counts in the
  low-degree sweep above \\
Textual audit of the $d\le6$ reconstruction &
  \art{results/steinberger/degree56-verification.md} (quote-by-quote against the
  pinned paper source) \\
Strong-encoding negative control: full pool \textsc{infeasible}, $y^\top b=-175$ &
  \art{tools/lp\_wheel\_level.py -{}-control},
  \art{results/steinberger/wheel\_level\_control\_stats.json} \\
$77{,}185$ leaves recovered from $833$ aborted jobs &
  \art{tools/lp\_rerun\_failed\_wheels.sh},
  \art{results/p3/runs/steinberger-s1/work/wheels/zero\_ndebug/},
  \art{results/steinberger/rerun\_ndebug.log} \\
S1: $833$ of $16{,}157$ jobs abort; $780$ \texttt{C\_positive} + $53$
  \texttt{degree\_range}; $530/250/51/2$ by hub degree &
  \art{results/p3/runs/steinberger-s1/failures.json};
  job counts from the \texttt{.done}/\texttt{.FAILED} markers under
  \art{results/p3/runs/steinberger-s1/log/} \\
Full narrative, scope and adversarial pass &
  \art{results/steinberger/LP-SCHEMA-RESULT.md} \\
\emph{Narrative-only} (recorded in prose in
  \art{LP-SCHEMA-RESULT.md}, no JSON): the $2{,}145$ monotonicity parent/child pairs
  and their $8$ failures under negative amounts; the six solver configurations and the
  $0.082$ uniform slack; the degree-$5$-alone / degree-$6$-alone ablations; the
  $\approx$14\,min / $\approx$2\,GB cost of \texttt{-{}-all}
  & re-derivable only by re-running \art{tools/lp\_discharge.py}; we flag these as
  the claims in \S\ref{sec:lp} not backed by a stored machine-readable artifact \\
\midrule
\multicolumn{2}{l}{\emph{\S\ref{sec:fr}: the $f(r)$ theorem}}\\[2pt]
Per-$(r,n)$ sweep counts (Table~\ref{tab:frsweep}) &
  \art{results/theorem/fr\_table/manifest\_r\{8..13\}.json} (r13 additionally
  \texttt{\_s\{0..7\}of8}); every entry has \texttt{d\_reducible\_count: 0} \\
Per-configuration records (adjacency, $a$, closure trace, legality flag) &
  \art{results/theorem/fr\_table/configs\_r*\_n*.jsonl} \\
Cross-check vs.\ the compiled 1995 oracle: $70$ sampled, $0$ mismatches,
  $6$ inconclusive &
  \art{results/theorem/fr\_table/crosscheck.json};
  \art{tools/fr\_table\_crosscheck.py} \\
Third-implementation adjudication of the $6$ &
  \art{src/fourcolor/brute\_force.py}, \art{tests/test\_brute\_force.py} \\
Tightness witnesses at $r=8..12$ &
  \art{results/theorem/fr\_table/tightness\_witnesses.json} \\
Tightness at $r=13$ ($k=8$, $a=6{,}954$) and the per-ring minimum-$k$ table &
  the labelled corpus via \art{src/fourcolor/lemma\_corpus.py} \\
C-vs-D witnesses at rings $8$ and $11$ &
  \art{third\_party/arxiv-1401.6481/src/anc/unavoidable.conf} idents
  \texttt{2.126}, \texttt{126.7566}, \texttt{128.136}, \texttt{7562.126} \\
Mass-law gap table (\S\ref{sec:gap}) &
  recomputed from \art{results/theorem/fr\_table/configs\_r*.jsonl}
  (max $a$ over legal below-threshold configs) against the corpus minima \\
Enumeration completeness &
  \art{third\_party/plantri/plantri55/} with sha256 of binary and source;
  legality by \texttt{fourcolor.mutate.is\_legal\_configuration} \\
Statement and protocol & \art{results/theorem/fr\_table/THEOREM.md} \\
\midrule
\multicolumn{2}{l}{\emph{\S\ref{sec:verify}: corpus verification}}\\[2pt]
RSST $633$: $633/633$ agree; $249$ D-reducible, $384$ C-reducible &
  \art{results/differential-unavoidable-r14/report.jsonl};
  \art{results/oracle-rsst-633/stdout.txt} \\
Steinberger $2{,}822$: $2{,}822/2{,}822$ agree, all D-reducible, all $b=0$ &
  \art{results/differential-U\_2822-r16/report.jsonl},
  \art{results/differential-stein-log.txt} \\
The Steinberger C oracle aborts at config $783$ &
  \art{results/oracle-stein-2822/stdout.txt} (ends \texttt{exit=134});
  crashing input \art{build/crash783.conf} \\
$7{,}697$ of the pruned pool verified, $73$ at rings $17$--$18$ unverified &
  \art{results/p3/pool-verify/SUMMARY.json}, \art{shard\_r*.jsonl} \\
2026 pipeline reproduced end to end, all counts matching &
  \art{results/nl4ct-reproduction-summary.txt};
  \art{third\_party/computer-checks/log/} \\
Differential on all $5{,}692{,}937$ wheel candidates, $0$ FP / $0$ FN &
  \art{results/nl4ct-differential/report\_d\{7..11\}.json}
  (\texttt{exact\_match: true} in each) \\
RSST 1995 discharging replay, $5/5$ presentation files &
  \art{results/discharge-rsst/replay.txt} \\
Wheel-level coverage: $85$ distinct blockers; cartwheel-level: $7{,}422$ over
  $2{,}253{,}185$ events &
  \art{results/p2-coverage/report.md},
  \art{results/p2-coverage/wheel\_level\_stats.json},
  \art{results/p2-coverage/cartwheel\_level\_stats.json} \\
$430$ zero-usage configurations &
  \art{results/p3/usage\_table\_v2\_summary.json},
  \art{results/p3/zero4\_configs.txt} (430 entries) \\
Pruned pool of $7{,}770$ passes the whole pipeline with identical counts &
  \art{results/p3/runs/batch1b/} --- \art{VERDICT.txt} (\textsc{pass}),
  \art{counts.json}, \art{checksums.sha256} \\
The failed $673$-config prune that revealed the fourth role &
  \art{results/p3/runs/batch1/VERDICT.txt} (\textsc{fail}, rc 134) \\
\midrule
\multicolumn{2}{l}{\emph{\S\ref{sec:mass}: the coloring-mass law}}\\[2pt]
Corpus of $59{,}142$, signature \texttt{a2b1feba331bf2962b2b}, $13{,}169$ D-reducible &
  \art{src/fourcolor/lemma\_corpus.py} (recomputed live);
  \art{results/mass-law/lemma\_log.jsonl} \\
Wheel closed form ($k=1$, support 6) & \textsc{corpus} receipt
  \texttt{b2fb2220ca3e6953}; proof in \S\ref{sec:mass}; four-way adversarial
  re-verification in \art{results/mass-law/PROOF-SHARP-CAP.md} \S6.3 \\
$a \le 11\cdot2^{r+k-7}$; $\min$-over-$h$ form tight only at the six wheels &
  receipts \texttt{e8beefcd5e56f5c5}, \texttt{4db503ddffcef7ee};
  \texttt{9977aea202daca71} (the killed $5\cdot2^{r+k-6}$) \\
Star-order claim verified on all $431{,}159$ (record, interior vertex) pairs &
  \art{tools/star\_order\_check.py},
  \art{results/mass-law/star-order-verification.md} \\
Exactly two common neighbours on all $610{,}846$ interior--interior edges &
  receipt \texttt{663fe16bb1a8ea3c} \\
(H2) false: arc histogram, $\max\gamma_{\rm prov}=1.937082$ &
  \art{tools/arc\_profile.py}, \art{results/mass-law/arc-profile.md} \\
Obstruction theorem: $P(S_0,4)=1464$, $a=55$ at \texttt{gen-r8-n10-2};
  $P$-form killed on $1{,}960$ of $59{,}142$ &
  \art{src/fourcolor/count4.py}, \art{tools/p4\_measure.py},
  \art{results/mass-law/p4-summary.md},
  \art{results/mass-law/p4\_measurements.jsonl} \\
Per-step factor table ($b_i=3$ mean $1.3155$, max $1.6479$) &
  \art{results/mass-law/star-order-verification.md} Part 4 \\
Lattice disks: $P$-form fails by $1.1\times10^4$; $a$ satisfies the sharp cap
  with improving margin &
  \art{tools/lattice\_disk\_entropy.py},
  \art{tools/lattice\_disk\_phi.py},
  \art{results/mass-law/lattice-disk-\{entropy,phi\}.md} \\
Threshold frontier, $\beta=2.4$ residual band, (T1) optimal to $0.08\%$ &
  \art{results/mass-law/threshold-candidates.md} (75 receipts) \\
Implied $f(r)$ tables &
  \art{results/mass-law/threshold-candidates.md} \S4,
  \art{results/mass-law/PROOF-SHARP-CAP.md} \S7 \\
Weak cap and Bridge Lemma & \art{results/theorem/mass-law/PROOF-CAP.md} \\
\midrule
\multicolumn{2}{l}{\emph{\S\ref{sec:ml} and \S\ref{sec:negative}}}\\[2pt]
D1 training and probe ladder &
  \art{results/d1/train\_metrics\_full\_v1.json},
  \art{results/d1/interp/report.md} \\
Boundary result, extraction, seed replication, error-set digestion &
  \art{results/d1v2/v2run/interrogation.md},
  \art{results/d1v2/extraction/report.md},
  \art{results/d1v2/replication/report.md},
  \art{results/d1v2/digestion/cases.md} \\
Zeta-map replication (companion repository) &
  companion repository \texttt{zeta\_map\_interp}: \texttt{results/train-metrics.json},
  \texttt{results/generalization.json},
  \texttt{results/interp/\{probes,ablation,attention-stats\}.json} \\
Mined candidate rules, coverage, cross-validation &
  \art{results/theorem/candidates.md},
  \art{results/theorem/mining\_results.json};
  \art{tools/mine\_sound\_rules.py} \\
Mutation stress test: $339{,}600$ mutants, $1{,}952$ labelled, $236$ FPs,
  $215$ oracle-confirmed &
  \art{results/theorem/stress\_test.json},
  \art{results/theorem/stress\_test.md} (every FP carries full adjacency and
  feature vector); \art{tools/stress\_test\_theorem.py} \\
Killed threshold forms & \textsc{corpus} receipts in
  \art{results/mass-law/lemma\_log.jsonl}; tabulated in
  \art{results/mass-law/threshold-candidates.md} \\
\bottomrule
\end{longtable}
}

%======================================================================
\section{The Farkas certificate}
\label{sec:certificate}
%======================================================================

Table~\ref{tab:cert} gives the $19$-row support of the certificate for
Theorem~\ref{thm:lp}. Exact totals: $\min_k (y^\top A)_k = 0$ (need $\ge 0$),
$y^\top b = -30$ (need $<0$), all $y_i > 0$. Every row is non-strict, so the
certificate proves infeasibility over $\mathbb{R}^{84}_{\ge0}$ and not merely over
the integers. The greedy-reduced $12$-row irreducible core is marked in the
\textsc{iis} column, with its own multipliers $y'$; it has $y'^\top b = -20$ and is
irreducible in the strict sense that deleting any single one of its $12$ rows
restores feasibility (verified for all $12$). Machine-readable:
\art{results/steinberger/lp\_main\_certificate.json} and
\art{results/steinberger/lp\_main\_certificate\_iis.json}; verify with
\art{tools/verify\_farkas.py}.

\begin{table}[h]
\centering\small
\caption{The $19$-row Farkas support ($y$, $y^\top b = -30$) and the $12$-row
irreducible core ($y'$, $y'^\top b = -20$; a dash marks a row absent from the core).
$C(x_0)$ is the leaf charge at the published amounts; a positive value marks a row
the published amounts violate. All four ring-$\le14$-only rows survive into the core.}
\label{tab:cert}
\begin{tabular}{rrrlrrll}
\toprule
$y_i$ & $y'_i$ & $d$ & spoke degrees & rhs & $C(x_0)$ & source & origin \\
\midrule
19 & 75 & 7 & 5,5,8,5,6,5,7 & 10 & 0 & \texttt{d7\_1155\_1}   & full $\cap$ r14 \\
21 & --- & 7 & 5,5,5,8,5,5,7 & 10 & 0 & \texttt{d7\_166\_0}    & full $\cap$ r14 \\
17 & 75 & 7 & 5,5,5,8,5,5,7 & 10 & 0 & \texttt{d7\_166\_1}    & full $\cap$ r14 \\
14 & --- & 7 & 5,6,5,6,7,5,8 & 10 & 0 & \texttt{d7\_1782\_2}   & full $\cap$ r14 \\
12 & --- & 7 & 5,6,5,8,5,7,6 & 10 & 0 & \texttt{d7\_1932\_3}   & full $\cap$ r14 \\
4  & 25 & 7 & 5,5,7,5,5,7,7 & 10 & 0 & \texttt{d7\_710\_10}   & full $\cap$ r14 \\
4  & --- & 7 & 5,5,7,5,5,7,7 & 10 & 0 & \texttt{d7\_710\_14}   & full $\cap$ r14 \\
21 & 75 & 7 & 5,5,7,5,6,5,8 & 10 & 0 & \texttt{d7\_722\_3}    & full $\cap$ r14 \\
9  & --- & 7 & 5,5,7,5,7,5,7 & 10 & 0 & \texttt{d7\_740\_16}   & full $\cap$ r14 \\
8  & --- & 7 & 5,5,7,5,7,5,7 & 10 & 0 & \texttt{d7\_740\_18}   & full $\cap$ r14 \\
13 & 55 & 7 & 5,5,7,5,7,5,7 & 10 & 0 & \texttt{d7\_740\_19}   & full $\cap$ r14 \\
20 & 60 & 7 & 5,7,5,7,6,7,7 & 10 & \textbf{2} & \texttt{d7\_3358\_2133} & \textbf{ring$\le$14 only} \\
6  & 15 & 7 & 5,7,5,7,6,7,7 & 10 & \textbf{2} & \texttt{d7\_3358\_2134} & \textbf{ring$\le$14 only} \\
30 & 55 & 7 & 5,7,5,7,7,6,7 & 10 & \textbf{2} & \texttt{d7\_3372\_1881} & \textbf{ring$\le$14 only} \\
4  & 25 & 7 & 5,7,5,7,7,7,7 & 10 & \textbf{3} & \texttt{d7\_3376\_1132} & \textbf{ring$\le$14 only} \\
2  & --- & 5 & 5,5,6,7,7     & $-10$ & $-3$ & \texttt{d5-55677} & low-degree row \\
43 & 75 & 5 & 5,5,7,6,7     & $-10$ & $-2$ & \texttt{d5-55767} & low-degree row \\
75 & 168 & 5 & 5,7,5,7,7     & $-10$ & 0    & \texttt{d5-57577} & low-degree row \\
85 & 219 & 5 & 5,7,6,7,7     & $-10$ & 0    & \texttt{d5-57677} & low-degree row \\
\bottomrule
\end{tabular}
\end{table}

\end{document}
